\documentclass[11pt]{article}
\usepackage[a4paper,margin=1in]{geometry}
\usepackage[T1]{fontenc}
\usepackage[utf8]{inputenc}
\usepackage{lmodern}
\usepackage{amsmath,amssymb,amsthm,mathtools}
\usepackage{enumitem}
\usepackage[hidelinks]{hyperref}
\usepackage{microtype}
\setlist{nosep}
\sloppy

\title{Atlas of Controlled Transfer Packages}
\author{Luca Blanchi}
\date{June 2026}

\begin{document}

\begin{abstract}
This addendum collects an operational atlas of transfer packages for Presentation Theory. It is not presented as a sequence of closed new theorems, but as a reusable map of contexts, descriptions, observables, verification data, fibres, scales, and overhead functions. Each network indicates which objects should be compared, which observables should be pulled back or pushed forward, which verification data must remain controllable, and which fibre geometry measures the loss of information. The purpose is to provide a complete repertoire of templates from which later work can extract transfer theorems with explicit hypotheses and budgets.
\end{abstract}

\section{Status and Reading Guide}
The atlas should be read as a foundational support document. An arrow between two theories does not automatically assert the existence of an equivalence or of a complete transfer theorem. It records instead the type of compiler, observational pullback, verification datum, fibre, and overhead that would have to be controlled in order to obtain a theorem. In this sense the central object is not an informal translation between domains, but the controlled package of access to information.

For each network one should distinguish three levels. The first is descriptive: it identifies contexts that are already naturally connected in mathematical practice. The second is formal: it specifies which maps between descriptions, observables, verification data, and fibres should enter the package. The third is demonstrative: it requires effective estimates for the overheads and suitable stability conditions. Only the third level produces an autonomous theorem; the first two provide a grammar for formulating it without losing the budgets.

The terminology used here follows the current version of the theory: observables are the functions, invariants, or procedures that measure an object through a declared channel; observable profiles are the finite or partial images produced by those channels; verification data are the controllable pieces of information that justify a property in the source or in the target; fibres are classes of objects indistinguishable with respect to a fixed profile. Overhead functions measure how descriptive size, observational cost, verification cost, and fibre-navigation cost grow under transfer.

\section{Common Grammar of the Atlas}

\subsection{Presentational Contexts}

A node of the atlas is not merely a theory (T), but a theory equipped with an access structure:

\[
\mathfrak P=
(\mathcal X,\Gamma,\mathfrak O,\Sigma,\mathfrak V,\mathfrak F,\mathfrak M).
\]

Where:

\[
\mathcal X
\]

is the class of objects.

\[
\Gamma=(\mathcal D,\rho,\kappa)
\]

is the presentation system: descriptions, realization, and cost.

\[
\mathfrak O
\]

is the system of observables.

\[
\Sigma
\]

is the operational structure: sums, products, compositions, limits, gluing, mutations, evolutions.

\[
\mathfrak V
\]

is the system of verification data.

\[
\mathfrak F
\]

is the geometry of fibres.

\[
\mathfrak M
\]

is the possible structure of measure, density, probability, or distribution on the objects.

The complexity of an object is not an absolute property. It is relative to:

\[
\text{how the object is described}
\]

\[
\text{how the object is observed}
\]

\[
\text{how the object is verified}
\]

\[
\text{how one moves among equivalent descriptions.}
\]



\subsection{Transfer package}

A transfer package from a context \(\mathfrak P_A\) to a context \(\mathfrak P_B\) is denoted by:

\[
\mathfrak T:A\longrightarrow B.
\]

The ideal form contains:

\[
(F,\widehat F,F^\ast,\Sigma,V,\Phi,\mu).
\]

Where:

\[
F:\mathcal X_A\to\mathcal X_B
\]

transports objects.

\[
\widehat F:\mathcal D_A\to\mathcal D_B
\]

compiles descriptions.

\[
F^\ast:\mathfrak O_B\to\mathfrak O_A
\]

pulls back observables.

\[
\Sigma
\]

controls the behavior with respect to operations.

\[
V
\]

transports verification data.

\[
\Phi
\]

controls fibres, moves, and bridge profiles.

\[
\mu
\]

transports measures, densities, distributions, or probabilistic limits.

The transfer package carries an overhead vector:

\[
\mathbf f_{\mathfrak T}=
(f_{\mathrm{pres}},
f_{\mathrm{obs}},
f_{\mathrm{op}},
f_{\mathrm{ver}},
f_{\mathrm{fib}},
f_{\mathrm{meas}},
f_{\mathrm{lim}},
f_{\mathrm{reg}},
f_{\mathrm{scale}}).
\]

There is not a single cost. There are many costs.



\subsection{Basic Resources}

We use the following resources:

\[
P=\text{presentational cost}
\]

\[
O=\text{observational cost}
\]

\[
V=\text{verification cost}
\]

\[
B=\text{cost of navigation in the fibre}
\]

\[
M=\text{metric, probabilistic, or measure-theoretic cost}
\]

\[
R=\text{regularity}
\]

\[
S=\text{scale}
\]

\[
\varepsilon=\text{precision}
\]

\[
T=\text{observation time}
\]

\[
\Lambda=\text{frequency or spectral cutoff}
\]

\[
d=\text{degree}
\]

\[
q=\text{quantifier rank}
\]

\[
n=\text{dimension, size, number of samples, or number of states.}
\]



\subsection{General Principles}

\subsection*{Principle 1: one transfers access, not the object}

\[
\boxed{
\text{a map transports objects; a transfer package transports access to the objects.}
}
\]

An ordinary reduction says:

\[
x\mapsto F(x).
\]

A transfer package says:

\[
\text{descriptions of }x
\mapsto
\text{descriptions of }F(x),
\]

\[
\text{observables on }F(x)
\mapsto
\text{observables on }x,
\]

\[
\text{verification data in the target}
\mapsto
\text{verification data in the source},
\]

\[
\text{fibre of the target}
\leftrightarrow
\text{fibre of the source.}
\]



\subsection*{Principle 2: every bounded observation induces fibres}

Given an observable, or a family of bounded observables:

\[
\mathcal O_{\le r},
\]

one obtains a partition:

\[
x\sim_r y
\quad\Longleftrightarrow\quad
O(x)=O(y)
\text{ for every } O\in\mathcal O_{\le r}.
\]

The hidden complexity lives in the classes:

\[
[x]_r.
\]

Therefore:

\[
\boxed{
\text{the geometry of the fibre is the geometry of forgotten information.}
}
\]



\subsection*{Principle 3: transfer packages compose}

If:

\[
A\xrightarrow{\mathfrak T}B\xrightarrow{\mathfrak S}C,
\]

then:

\[
A\xrightarrow{\mathfrak S\circ\mathfrak T}C.
\]

The overheads compose:

\[
f_{\mathfrak S\circ\mathfrak T}
\le
f_{\mathfrak S}\circ f_{\mathfrak T}.
\]

This produces \textbf{transfer graphs}.



\subsection*{Principle 4: Morita theory with budgets}

Two theories may be considered equivalent, not absolutely, but up to a declared class of overhead functions \(\mathcal H\):

\[
A\simeq_{\mathcal H}B.
\]

This means that there are packages in both directions:

\[
A\to B,
\qquad
B\to A,
\]

whose compositions are equivalent to the identity with controlled overhead.

This is a \textbf{quantitative Morita theory}.



\subsection*{Principle 5: every theory has many observable profiles}

Every theory produces many observable profiles:

\[
\text{spectra}
\]

\[
\text{moments}
\]

\[
\text{entropies}
\]

\[
\text{cohomologies}
\]

\[
\text{finite quotients}
\]

\[
\text{truncations}
\]

\[
\text{samples}
\]

\[
\text{invariants}
\]

\[
\text{categories}
\]

\[
\text{operators}
\]

\[
\text{kernel}
\]

\[
\text{tensors}.
\]

The task of the atlas is to measure the cost of passing from one observable profile to another.



\section{The Main Transversal Hubs}

The atlas is large, but many edges pass through a small number of universal hubs.



\subsection{Hub of Observables}

\[
\mathcal O(\mathcal X)
\]

is the algebra, category, or system of observables of a theory.

Every object gives an evaluation:

\[
x\mapsto \operatorname{ev}_x:\mathcal O(\mathcal X)\to A.
\]

Every map:

\[
F:\mathcal X\to\mathcal Y
\]

induces a pullback:

\[
F^\ast:\mathcal O(\mathcal Y)\to\mathcal O(\mathcal X).
\]

This is the contravariant core of the transfer theory.



\subsection{Hub of Fibres}

Every observational map:

\[
F:\mathcal X\to\mathcal Y
\]

produces fibres:

\[
F^{-1}(y).
\]

Fibres should be studied with:

\[
\text{moves}
\]

\[
\text{bridge profiles}
\]

\[
\text{diameters}
\]

\[
\text{mixing}
\]

\[
\text{singularities}
\]

\[
\text{moduli}
\]

\[
\text{internal probabilities}
\]

\[
\text{navigation algorithms}.
\]

Universal examples:

\[
\text{functions with the same low Fourier coefficients}
\]

\[
\text{domains with the same first eigenvalues}
\]

\[
\text{metrics with the same truncated heat trace}
\]

\[
\text{PDEs with the same input-output data}
\]

\[
\text{neural networks representing the same function}
\]

\[
\text{distributions with same moments}
\]

\[
\text{graphs with the same small subgraphs}
\]

\[
\text{strutture with same theory logic up to a quantifier rank }q.
\]



\subsection{Hub of Operators}

Many theories are transported into operators:

\[
X\mapsto T_X.
\]

Examples:

\[
M\mapsto \Delta_M
\]

\[
G\mapsto L_G
\]

\[
T:X\to X\mapsto U_T
\]

\[
\text{Markov process}\mapsto L
\]

\[
\text{kernel }K(x,y)\mapsto T_K
\]

\[
\text{PDE}\mapsto P(x,D).
\]

Then:

\[
T\mapsto \operatorname{Spec}(T),
\]

\[
T\mapsto e^{tT},
\]

\[
T\mapsto (z-T)^{-1},
\]

\[
T\mapsto \operatorname{Tr}(e^{-tT}).
\]



\subsection{Hub of Kernels}

\[
K(x,y)
\]

appears in:

\[
\text{PDE}
\]

\[
\text{machine learning}
\]

\[
\text{operator theory}
\]

\[
\text{probability}
\]

\[
\text{graph limits}
\]

\[
\text{quantum mechanics}
\]

\[
\text{integral geometry}.
\]

Scheme:

\[
\begin{gathered}
\text{object} \to \text{kernel} \to \text{operator} \to \text{spectrum} \\
\to \text{features}.
\end{gathered}
\]



\subsection{Spectral Hub}

\[
\operatorname{Spec}
\]

connects:

\[
\text{graphs}
\]

\[
\text{varieties}
\]

\[
\text{operators}
\]

\[
\text{dynamics}
\]

\[
\text{quantum mechanics}
\]

\[
\text{automorphic forms}
\]

\[
\text{zeta functions}
\]

\[
\text{trace formulas}.
\]

Scheme:

\[
\begin{gathered}
\text{object} \to \text{operator} \to \text{spectrum} \to \text{zeta/trace} \\
\to \text{arithmetic or geometry}.
\end{gathered}
\]



\subsection{Hub of Moments}

\[
m_\alpha=\int x^\alpha,d\mu.
\]

Moments connect:

\[
\text{probability}
\]

\[
\text{statistics}
\]

\[
\text{SDP}
\]

\[
\text{operator algebras}
\]

\[
\text{quantum correlations}
\]

\[
\text{inverse problems}.
\]

Scheme:

\[
\begin{gathered}
\text{object} \to \text{moments} \to \text{semidefinite relaxation} \to \text{verification datum}.
\end{gathered}
\]



\subsection{Entropic Hub}

\[
H,\quad S,\quad \operatorname{Ent},\quad h_{\mathrm{top}},\quad h_\mu.
\]

Connects:

\[
\text{probability}
\]

\[
\text{dynamics}
\]

\[
\text{compression}
\]

\[
\text{statistical mechanics}
\]

\[
\text{optimal transport}
\]

\[
\text{learning theory}
\]

\[
\text{quantum information}.
\]

Scheme:

\[
\begin{gathered}
\text{object} \to \text{distribution} \to \text{entropy} \to \text{concentration/compression/mixing}.
\end{gathered}
\]



\subsection{Sheaf Hub}

\[
\mathcal F.
\]

Connects:

\[
\text{topology}
\]

\[
\text{PDE microlocali}
\]

\[
\text{database}
\]

\[
\text{CSP}
\]

\[
\text{distributed systems}
\]

\[
\text{algebraic geometry}
\]

\[
\text{persistence}.
\]

Scheme:

\[
\begin{gathered}
\text{local data} \to \text{sheaf} \to \text{cohomology} \to \text{obstruction}.
\end{gathered}
\]



\subsection{Tensorial Hub}

\[
T\in V_1\otimes\cdots\otimes V_k.
\]

Connects:

\[
\text{algebra}
\]

\[
\text{quantum information}
\]

\[
\text{PDE kernels}
\]

\[
\text{statistics}
\]

\[
\text{machine learning}
\]

\[
\text{complexity theory}
\]

\[
\text{cohomology}.
\]

Scheme:

\[
\begin{gathered}
\text{struttura multilineare} \to \text{tensore} \to \text{rank} \to \text{lower bound}.
\end{gathered}
\]



\subsection{Hub of Convex Verification Data}

Includes:

\[
\text{SOS}
\]

\[
\text{SDP}
\]

\[
\text{dual verification data}
\]

\[
\text{Lyapunov functions}
\]

\[
\text{entropy verification data}
\]

\[
\text{barrier verification data}
\]

\[
\text{moment relaxations}.
\]

Scheme:

\[
\text{truth}
\to
\text{verification datum}
\to
\text{verifiable bound}.
\]



\subsection{Hub of Approximations}

All of these are variants of the same scheme:

\[
X\mapsto X_{\le r}.
\]

Examples:

\[
\text{Taylor jets}
\]

\[
\text{Fourier cutoffs}
\]

\[
\text{finite quotients}
\]

\[
\text{mesh refinements}
\]

\[
\text{moment truncations}
\]

\[
\text{logical quantifier-rank truncations}
\]

\[
\text{spectral truncations}
\]

\[
\text{prefixes of words}
\]

\[
\text{finite samples}.
\]



\subsection{Hub of Dualities}

Dualities are bidirectional transfer packages.

Examples:

\[
\text{compact spaces}
\leftrightarrow
\text{commutative }C^\ast\text{-algebras}
\]

\[
\text{locally compact abelian groups}
\leftrightarrow
\text{Pontryagin duals}
\]

\[
\text{corpi convessi}
\leftrightarrow
\text{support functions}
\]

\[
\text{measures}
\leftrightarrow
\text{moments/characteristic functions}
\]

\[
\text{D-modules}
\leftrightarrow
\text{perverse sheaves}
\]

\[
\text{symplectic}
\leftrightarrow
\text{mirror algebraic geometry}
\]

\[
\text{syntax}
\leftrightarrow
\text{semantics}
\]

\[
\text{automorfo}
\leftrightarrow
\text{Galois}.
\]

The presentational question is not only whether the duality exists, but:

\[
\text{how much do descriptions, observables, and verification data grow across the duality?}
\]



\section{Atlas of Networks}

We now pass to the networks themselves.

Each network is a graph of theories. Each edge is a possible transfer package.



\section{A. Foundational, Logical, and Syntactic Networks}



\subsection{A1. Network syntax--semantics--CSP--database--proof complexity}

Nodes:

\[
\text{equational theories}
\]

\[
\text{universal algebras}
\]

\[
\text{Lawvere theories}
\]

\[
\text{operadi}
\]

\[
\text{monadi}
\]

\[
\text{term rewriting}
\]

\[
\text{CSP}
\]

\[
\text{database theory}
\]

\[
\text{finite moof the theory}
\]

\[
\text{proof complexity}.
\]

Main edges:

\[
\text{presentation equazionale}
\to
\text{category of models}
\]

\[
\text{equational theory}
\to
\text{CSP template}
\]

\[
\text{database schema}
\leftrightarrow
\text{finite-limit theory}
\]

\[
\text{term rewriting}
\to
\text{proof complexity}.
\]

Resources:

\[
P=\text{number of symbols, arity, term length}
\]

\[
O=\text{quantifier rank, width, treewidth}
\]

\[
V=\text{derivation length, degree, width}
\]

\[
B=\text{distance between equivalent presentations}.
\]

Central idea:

\[
\text{query of database}
\Longleftrightarrow
\text{observable}
\]

\[
\text{proof derivation}
\Longleftrightarrow
\text{bridge in the fibra of the presentations equivalent}.
\]



\subsection{A2. Network types--proofs--categories--homotopy--formalization}

Nodes:

\[
\lambda\text{-calculus}
\]

\[
\text{type theory}
\]

\[
\text{dependent type theory}
\]

\[
\text{proof assistants}
\]

\[
\text{categories}
\]

\[
\text{toposes}
\]

\[
\text{homotopy type theory}
\]

\[
\infty\text{-categories}
\]

\[
\text{rewriting}
\]

\[
\text{normalization}
\]

\[
\text{proof data}.
\]

Edges:

\[
\text{proof}
\to
\text{term}
\]

\[
\text{type theory}
\to
\text{category}
\]

\[
\text{category}
\to
\text{internal language}
\]

\[
\text{homotopy type}
\to
\text{space/simplicial object}
\]

\[
\text{formal proof}
\to
\text{checkable proof data}.
\]

Resources:

\[
P=\text{context dimension, universe levels, costruttori}
\]

\[
O=\text{depth of types, truncation level}
\]

\[
V=\text{type-checking time, proof-term length}
\]

\[
B=\text{distance between equivalent proofs, normalization length}.
\]

Central question:

\[
\text{can a mathematically simple result be formally expensive?}
\]



\subsection{A3. Network continuous logic--functional analysis--probability}

Nodes:

\[
\text{continuousus logic}
\]

\[
\text{metric structures}
\]

\[
\text{Banach spaces}
\]

\[
C^\ast\text{-algebras}
\]

\[
\text{probability algebras}
\]

\[
\text{tracial von Neumann algebras}
\]

\[
\text{ultraproducts metrici}
\]

\[
\text{stability theory}
\]

\[
\text{definability}.
\]

Edges:

\[
X\mapsto \operatorname{Th}_{\le q}(X)
\]

\[
(X_n)\mapsto\prod_{\mathcal U}X_n
\]

\[
A\mapsto \text{tracial moments}
\]

\[
\text{Banach space}\mapsto \text{metric theory}.
\]

Resources:

\[
O=\text{quantifier rank, modulus of continuity, precision}
\]

\[
V=\text{elementary approximation verification data}
\]

\[
B=\text{fibres of indistinguishable structures up to a rank }q.
\]



\subsection{A4. Network o-minimality--PDE tame--real geometry--algorithms}

Nodes:

\[
\text{o-minimal structures}
\]

\[
\text{semialgebraic sets}
\]

\[
\text{subanalytic sets}
\]

\[
\text{Pfaffian functions}
\]

\[
\text{definable manifolds}
\]

\[
\text{tame topology}
\]

\[
\text{real analytic geometry}
\]

\[
\text{optimization}
\]

\[
\text{differential equations definable}.
\]

Edges:

\[
f\mapsto \operatorname{Graph}(f)
\]

\[
\text{definable set}\mapsto \text{Betti numbers, cell decomposition}
\]

\[
\text{ODE definibile}\mapsto \text{flow definibile}
\]

\[
\text{optimization problem}\mapsto \text{critical point data}.
\]

Resources:

\[
P=\text{format definibile, degree, chain length}
\]

\[
O=\text{cell dimension, Betti complexity}
\]

\[
V=\text{stratification verification data, monotonicity}
\]

\[
B=\text{fibres with the same tame decomposition}.
\]

Idea:

\[
\text{analysis with controlled topological complexity}.
\]



\subsection{A5. Network finite precision--computable analysis--real complexity}

Nodes:

\[
\text{computable real numbers}
\]

\[
\text{represented spaces}
\]

\[
\text{Type-2 computability}
\]

\[
\text{computable analysis}
\]

\[
\text{Weihrauch degrees}
\]

\[
\text{numerical algorithms}
\]

\[
\text{real RAM models}
\]

\[
\text{interval computation}
\]

\[
\text{constructive analysis}.
\]

Edges:

\[
x\in\mathbb R\mapsto \text{name }(q_n)
\]

\[
F:X\to Y\mapsto \text{realizer}
\]

\[
\text{analytic theorem}\mapsto \text{Weihrauch degree}
\]

\[
\text{numerical approximation}\mapsto \text{interval verification data}.
\]

Resources:

\[
O=2^{-n}\text{ precision}
\]

\[
V=\text{modulus of continuity, convergenza, enclosure}
\]

\[
B=\text{fibre of different names of the same object}.
\]



\subsection{A6. Network PDE--logic--proof mining--constructive analysis}

Nodes:

\[
\text{proof mining}
\]

\[
\text{constructive analysis}
\]

\[
\text{functional interpretation}
\]

\[
\text{fixed point theorems}
\]

\[
\text{ergodic theorems}
\]

\[
\text{PDE existence proofs}
\]

\[
\text{compactness arguments}
\]

\[
\text{rates of convergence}
\]

\[
\text{metastability}.
\]

Edges:

\[
\text{classical proof}\mapsto \text{effective bound}
\]

\[
\text{convergence theorem}\mapsto \text{metastability bound}
\]

\[
\text{compactness proof}\mapsto \text{modulus extraction}
\]

\[
\text{PDE existence}\mapsto \text{approximation scheme + bound}.
\]

Idea:

\[
\begin{gathered}
\text{qualitative theorem} \to \text{quantitative profile} \to \text{algorithm} \to \text{verification datum}.
\end{gathered}
\]



\section{B. Algebraic, Representational, and Homological Networks}



\subsection{B1. Network representations--moduli--invariants--deformations}

Nodes:

\[
\text{finitely presented groups}
\]

\[
\text{associative algebras}
\]

\[
\text{Lie algebras}
\]

\[
\text{quiver}
\]

\[
\text{rings}
\]

\[
\text{tensor categories}
\]

\[
\text{representation varieties}
\]

\[
\text{character varieties}
\]

\[
\text{GIT quotients}
\]

\[
\text{moduli stacks}
\]

\[
\text{deformation theory}.
\]

Edges:

\[
G\to \operatorname{Rep}_n(G)
\]

\[
A\to \operatorname{Mod}_n(A)
\]

\[
Q\to \operatorname{Rep}(Q,\mathbf d)
\]

\[
\mathcal C\to K_0(\mathcal C)
\]

\[
X\to \mathfrak g_X.
\]

Observables:

\[
\text{traces of words}
\]

\[
\text{characters}
\]

\[
\text{semi-invariants}
\]

\[
\text{stabilizers}
\]

\[
\text{orbit closures}
\]

\[
\text{deformation cohomology}.
\]

Fibres:

\[
\text{non-isomorphic objects with the same invariants}
\]

\[
\text{orbits with the same closure}
\]

\[
\text{moduli with the same decategorification}.
\]



\subsection{B2. Universal Tensorial Network}

Nodes:

\[
\text{finite-dimensional algebras}
\]

\[
\text{matrix multiplication}
\]

\[
\text{groups and group algebras}
\]

\[
\text{cohomology and cup product}
\]

\[
\text{multilinear forms}
\]

\[
\text{probability discrete}
\]

\[
\text{graphical models}
\]

\[
\text{quantum states}
\]

\[
\text{tensor networks}
\]

\[
\text{PDE kernels}
\]

\[
\text{complexity of bilinear maps}.
\]

Edges:

\[
A\mapsto T_A
\]

where \(T_A\) is the multiplication tensor.

\[
X\mapsto T_{H^\ast(X)}
\]

cup product in cohomology.

\[
p(x_1,\ldots,x_n)\mapsto T_p
\]

distribution come tensore.

\[
\psi\mapsto T_\psi
\]

stato quantistico multipartito.

\[
K(x,y)\mapsto T_K
\]

kernel discretizzato or approssimato.

Observables:

\[
\text{rank}
\]

\[
\text{border rank}
\]

\[
\text{slice rank}
\]

\[
\text{partition rank}
\]

\[
\text{flattening rank}
\]

\[
\text{bond dimension}
\]

\[
\text{entanglement entropy}.
\]

Idea:

\[
\text{every composizione bilineare or multilineare genera a transfer tensoriale.}
\]



\subsection{B3. Network linear algebra--matroids--codes--data structures--rank methods}

Nodes:

\[
\text{matrices}
\]

\[
\text{matroids}
\]

\[
\text{linear codes}
\]

\[
\text{linear data structures}
\]

\[
\text{communication matrices}
\]

\[
\text{incidence geometries}
\]

\[
\text{polynomial method}
\]

\[
\text{linear circuits}
\]

\[
\text{proof complexity algebrica}.
\]

Edges:

\[
\text{problem}\mapsto A_P
\]

matrice of incidenza, comunicazione or vincoli.

\[
A\mapsto M(A)
\]

matroid.

\[
A\mapsto C_A
\]

linear code.

\[
\text{data structure}\mapsto A=RM
\]

fattorizzazione sparsa.

\[
\text{polynomial proof}\mapsto \text{linear map of coefficients}.
\]

Resources:

\[
P=\text{dimension, sparsity, field}
\]

\[
O=\text{rank, local rank, support size}
\]

\[
V=\text{dependence verification data, dual codewords}
\]

\[
B=\text{distance between representations of the same matroid}.
\]



\subsection{B4. Network deformations--cohomology--obstructions--derived moduli}

Nodes:

\[
\text{algebraic varieties}
\]

\[
\text{schemes}
\]

\[
\text{complexes}
\]

\[
\text{dg-algebras}
\]

\[
\text{dg-Lie algebras}
\]

\[
L_\infty\text{-algebras}
\]

\[
\text{Hochschild cohomology}
\]

\[
\text{André--Quillen cohomology}
\]

\[
\text{deformation functors}
\]

\[
\text{derived stacks}
\]

\[
\text{obstruction theory}.
\]

Edges:

\[
X\mapsto \mathfrak g_X
\]

\[
\mathfrak g\mapsto MC(\mathfrak g)
\]

\[
MC(\mathfrak g)\mapsto H^\ast(\mathfrak g)
\]

\[
\text{moduli problem}\mapsto \text{derived moduli stack}.
\]

Resources:

\[
P=\text{dimension of the complesso, differentials}
\]

\[
O=\text{observed cohomological degree}
\]

\[
V=\text{nilpotent order of the verified deformation}
\]

\[
B=\text{gauge equivalence, homotopies, quasi-isomorphisms}.
\]

Idea:

\[
\text{nonlinear classification}
\to
\text{linear data plus higher obstructions}.
\]



\subsection{B5. Network homological algebra--functional analysis--derived categories}

Nodes:

\[
\text{complexes}
\]

\[
\text{derived categories}
\]

\[
\text{triangulated categories}
\]

\[
\text{dg-categories}
\]

\[
\text{Ext/Tor}
\]

\[
\text{Hochschild homology}
\]

\[
\text{cyclic homology}
\]

\[
\text{topological cyclic homology}
\]

\[
\text{operator K-theory}
\]

\[
\text{index theory}.
\]

Edges:

\[
A\mapsto D(A)
\]

\[
M,N\mapsto \operatorname{Ext}^\ast(M,N)
\]

\[
A\mapsto HH_\ast(A)
\]

\[
A\mapsto HC_\ast(A)
\]

\[
\text{smooth algebra}\mapsto \text{de Rham-type invariants}.
\]

Fibres:

\[
\text{categories with same invariants omologici troncati}.
\]



\section{C. Geometric, Topological, and Categorical Networks}



\subsection{C1. Network tropical--matroidal--polyhedral--optimization}

Nodes:

\[
\text{algebraic varieties}
\]

\[
\text{ideali}
\]

\[
\text{schemes}
\]

\[
\text{Newton polytopes}
\]

\[
\text{tropical varieties}
\]

\[
\text{matroids}
\]

\[
\text{oriented matroids}
\]

\[
\text{fans}
\]

\[
\text{cluster algebras}
\]

\[
\text{linear programming}
\]

\[
\text{integer programming}
\]

\[
\text{phylogenetic geometry}.
\]

Edges:

\[
I\mapsto \operatorname{Trop}(I)
\]

\[
I\mapsto \operatorname{in}_w(I)
\]

\[
\text{linear arrangement}\mapsto \text{matroid}
\]

\[
\text{varieties positiva}\mapsto \text{polytope}
\]

\[
\text{cluster algebra}\mapsto \text{fan combinatorics}.
\]

Resources:

\[
P=\text{degree, number of monomials, bit-size}
\]

\[
O=\text{dimension of the fan, number of cones, matroidal rank}
\]

\[
V=\text{tropical verification data, Groebner bases}
\]

\[
B=\text{distance of mutazione, distance in the Gröbner fan}.
\]



\subsection{C2. Network sheaf--cosheaf--local consistency--persistence--distributed systems}

Nodes:

\[
\text{sheaves}
\]

\[
\text{cosheaves}
\]

\[
\text{CSP locali}
\]

\[
\text{database joins}
\]

\[
\text{bundle theory}
\]

\[
\text{Cech cohomology}
\]

\[
\text{persistent homology}
\]

\[
\text{distributed computing}
\]

\[
\text{sensor networks}
\]

\[
\text{local-to-global proof systems}
\]

\[
\text{descent theory}.
\]

Edges:

\[
\text{CSP instance}\mapsto \mathcal F_{\text{solutions}}
\]

\[
\text{database}\mapsto \text{sheaf of compatible records}
\]

\[
\text{cover}\mapsto \check C^\bullet(\mathcal U,\mathcal F)
\]

\[
\text{filtered data}\mapsto \text{persistence module}
\]

\[
\text{distributed task}\mapsto \text{simplicial carrier/sheaf}.
\]

Idea:

\[
\begin{gathered}
\text{CSP width}\leftrightarrow\text{cohomological obstruction degree}\\
\leftrightarrow\text{database join complexity}\leftrightarrow\text{distributed task impossibility}.
\end{gathered}
\]



\subsection{C3. Network topology--groups--knots--TQFT--skein--mapping class groups}

Nodes:

\[
\text{complexes simpliciali}
\]

\[
\text{varieties PL}
\]

\[
\text{handle decompositions}
\]

\[
\text{knots and link}
\]

\[
\text{fundamental groups}
\]

\[
\text{quandles}
\]

\[
\text{skein modules}
\]

\[
\text{TQFT}
\]

\[
\text{mapping class groups}
\]

\[
\text{modular tensor categories}
\]

\[
\text{character varieties}.
\]

Edges:

\[
M\mapsto \pi_1(M)
\]

\[
M\mapsto C_\bullet(M)
\]

\[
K\mapsto \pi_1(S^3\setminus K)
\]

\[
K\mapsto Q(K)
\]

\[
M\mapsto Z_{\mathcal C}(M)
\]

\[
K\mapsto \text{skein element}.
\]

Observables:

\[
\text{homology}
\]

\[
\text{gruppo fondamentale}
\]

\[
\text{colorazioni quandle}
\]

\[
\text{quantum invariants up to level }r.
\]



\subsection{C4. Network persistent topology--functional analysis--signals--PDE}

Nodes:

\[
\text{filtered spaces}
\]

\[
\text{persistent homology}
\]

\[
\text{persistence modules}
\]

\[
\text{barcodes}
\]

\[
\text{sheaf persistence}
\]

\[
\text{signals}
\]

\[
\text{images}
\]

\[
\text{solutions of PDE}
\]

\[
\text{random fields}
\]

\[
\text{multiparameter persistence}.
\]

Edges:

\[
f:X\to\mathbb R\mapsto {X_{\le t}}_t
\]

\[
{X_t}\mapsto H_k(X_t)
\]

\[
\text{module}\mapsto \text{barcode}
\]

\[
u(x,t)\mapsto \text{persistence over level/time}
\]

\[
\text{random field}\mapsto \text{random barcode}.
\]

Central fibre:

\[
f\mapsto \operatorname{Barcode}(f)
\]

forgets a large part of the geometry.



\subsection{C5. Network symplectic geometry--Floer--mirror symmetry--elliptic PDE}

Nodes:

\[
\text{symplectic manifolds}
\]

\[
\text{Hamiltonian dynamics}
\]

\[
\text{Lagrangian submanifolds}
\]

\[
\text{pseudo-holomorphic curves}
\]

\[
\text{Floer homology}
\]

\[
\text{Fukaya categories}
\]

\[
\text{mirror symmetry}
\]

\[
\text{derived categories of coherent sheaves}
\]

\[
\text{Gromov--Witten theory}
\]

\[
\text{quantum cohomology}.
\]

Edges:

\[
(M,\omega)\mapsto \text{Hamiltonian flow}
\]

\[
L_0,L_1\mapsto CF^\ast(L_0,L_1)
\]

\[
J\mapsto \text{moduli of }J\text{-holomorphic curves}
\]

\[
X\mapsto QH^\ast(X)
\]

\[
\text{symplectic side}\mapsto \text{algebraic mirror}.
\]



\subsection{C6. Network Teichmüller--quasiconformal--complex dynamics}

Nodes:

\[
\text{Riemann surfaces}
\]

\[
\text{Teichmüller spaces}
\]

\[
\text{moduli of curves}
\]

\[
\text{quasiconformal maps}
\]

\[
\text{mapping class groups}
\]

\[
\text{measured foliations}
\]

\[
\text{geodesic flows}
\]

\[
\text{complex dynamics}
\]

\[
\text{flat surfaces}
\]

\[
\text{interval exchange transformations}.
\]

Edges:

\[
X\mapsto \mathcal T(X)
\]

\[
q\mapsto \text{horizontal/vertical foliations}
\]

\[
\text{flat surface}\mapsto \text{translation flow}
\]

\[
\text{mapping class}\mapsto \text{action on curves/homology}
\]

\[
\text{rational map}\mapsto \text{Julia set / lamination}.
\]



\subsection{C7. Network measure geometry--currents--varifold--fractals}

Nodes:

\[
\text{measures of Hausdorff}
\]

\[
\text{rectifiability}
\]

\[
\text{currents}
\]

\[
\text{varifolds}
\]

\[
\text{minimal surfaces}
\]

\[
\text{free boundary}
\]

\[
\text{fractals}
\]

\[
\text{multifractal analysis}
\]

\[
\text{geometric flows}
\]

\[
\text{GMT compactness}.
\]

Edges:

\[
S\mapsto \mathcal H^d\llcorner S
\]

\[
S\mapsto [S]
\]

\[
S\mapsto V_S
\]

\[
T\mapsto \partial T
\]

\[
\mu\mapsto \dim_H(\mu),\tau(q),f(\alpha).
\]

Observables:

\[
\text{scale}
\]

\[
\text{massa}
\]

\[
\text{flat norm}
\]

\[
\text{dimension}
\]

\[
\text{multifractal profile}.
\]



\subsection{C8. Network fractals--iterated dynamics--spectra--probability}

Nodes:

\[
\text{iterated function systems}
\]

\[
\text{self-similar sets}
\]

\[
\text{self-affine sets}
\]

\[
\text{fractal measures}
\]

\[
\text{analysis on fractals}
\]

\[
\text{spectral decimation}
\]

\[
\text{random fractals}
\]

\[
\text{multifractal formalism}
\]

\[
\text{substitution systems}.
\]

Edges:

\[
{S_i}\mapsto K
\]

\[
K\mapsto \mu_K
\]

\[
K\mapsto \Delta_K
\]

\[
\Delta_K\mapsto \operatorname{Spec}(\Delta_K)
\]

\[
\text{substitution}\mapsto \text{tiling dynamical system}.
\]



\section{D. Analytic Networks}

This is the large analytic expansion of the atlas.



\subsection{D1. Network distributions--Fourier--Sobolev--microlocal}

Nodes:

\[
\text{smooth functions}
\]

\[
L^p
\]

\[
\text{distributions}
\]

\[
\text{measures}
\]

\[
\text{trasformata of Fourier}
\]

\[
\text{Sobolev spaces}
\]

\[
\text{Besov spaces}
\]

\[
\text{wavelets}
\]

\[
\text{operators pseudodifferenziali}
\]

\[
\text{wavefront sets}
\]

\[
\text{microlocal sheaves}.
\]

Edges:

\[
u\mapsto \widehat u
\]

\[
u\mapsto |u|_{H^s}
\]

\[
u\mapsto WF(u)
\]

\[
P(x,D)\mapsto \sigma(P)
\]

\[
u\mapsto {\langle u,\psi_{j,k}\rangle}.
\]

Fibres:

\[
u\mapsto \widehat u|_{|\xi|\le \Lambda}
\]

forgets high frequencies.

\[
u\mapsto WF(u)
\]

forgets amplitudes and fine phases.

\[
P\mapsto \sigma_{\mathrm{prin}}(P)
\]

forgets subprincipal symbols.



\subsection{D2. Network elliptic PDE--regularity--geometry--topology}

Nodes:

\[
\text{operators ellittici}
\]

\[
\text{problemi al bordo}
\]

\[
\text{Sobolev spaces}
\]

\[
\text{Fredholm theory}
\]

\[
\text{index}
\]

\[
\text{cohomology}
\]

\[
\text{K-theory}
\]

\[
\text{Riemannian geometry}
\]

\[
\text{spectral theory}
\]

\[
\text{invariants topologici}.
\]

Edges:

\[
P\mapsto \sigma(P)
\]

\[
P\mapsto \operatorname{ind}(P)
\]

\[
Pu=f\mapsto u=P^{-1}f
\]

\[
(M,E,P)\mapsto [\sigma(P)]\in K(T^\ast M)
\]

\[
P\mapsto \ker P,\operatorname{coker}P.
\]

Idea:

\[
\text{PDE ellittica}
\to
\text{topology}
\to
\text{non-invertibility or existence verification data}.
\]



\subsection{D3. Network parabolic PDE--semigroups--probability--geometry}

Nodes:

\[
\text{heat equations}
\]

\[
\text{contractive semigroups}
\]

\[
\text{operators dissipativi}
\]

\[
\text{processi of Markov}
\]

\[
\text{Dirichlet forms}
\]

\[
\text{heat kernels}
\]

\[
\text{curvature-dimension}
\]

\[
\text{gradient flows}
\]

\[
\text{optimal transport}.
\]

Edges:

\[
L\mapsto e^{tL}
\]

\[
L\mapsto X_t
\]

\[
L\mapsto p_t(x,y)
\]

\[
p_t(x,y)\mapsto \operatorname{Tr}(e^{tL})
\]

\[
\partial_t u=Lu\mapsto \text{gradient flow}
\]

\[
\mathcal E(u,u)\mapsto \text{Dirichlet form}\mapsto \text{Markov process}.
\]



\subsection{D4. Network hyperbolic PDE--propagation--causality--scattering}

Nodes:

\[
\text{equazioni of the onde}
\]

\[
\text{operators iperbolici}
\]

\[
\text{geometry lorentziana}
\]

\[
\text{propagatori}
\]

\[
\text{fronti of onda}
\]

\[
\text{bicharacteristics}
\]

\[
\text{scattering}
\]

\[
\text{control theory}
\]

\[
\text{inverse boundary problems}.
\]

Edges:

\[
P\mapsto \operatorname{Char}(P)
\]

\[
\operatorname{Char}(P)\mapsto \text{flusso hamiltoniano}
\]

\[
u_0,u_1\mapsto u(t)
\]

\[
u(t)|_{\partial M}\mapsto \text{boundary observation}
\]

\[
(M,g)\mapsto \Lambda_g.
\]

Questions:

\[
\text{how much time is needed to see a singularity?}
\]

\[
\text{which angular aperture is needed to reconstruct a geometry?}
\]



\subsection{D5. Network dispersive--nonlinear--scattering--solitons}

Nodes:

\[
\text{Schrödinger nonlinear}
\]

\[
\text{wave maps}
\]

\[
\text{Klein--Gordon}
\]

\[
\text{KdV}
\]

\[
\text{Strichartz estimates}
\]

\[
\text{profile decompositions}
\]

\[
\text{concentration compactness}
\]

\[
\text{scattering states}
\]

\[
\text{solitons}.
\]

Edges:

\[
u_0\mapsto u(t)
\]

\[
u(t)\mapsto u_\pm
\]

\[
u_n\mapsto {\phi_j,\lambda_j,x_j,t_j}
\]

\[
u\mapsto \text{mass, energy, momentum}
\]

\[
\text{nonlinear PDE}\mapsto \text{renormalized effective dynamics}.
\]

Fibres:

\[
u_0\mapsto (M(u_0),E(u_0))
\]

has dynamically rich fibres.



\subsection{D6. Network harmonic analysis--singular operators--additive combinatorics}

Nodes:

\[
\text{Fourier analysis}
\]

\[
\text{Calderón--Zygmund operators}
\]

\[
\text{maximal functions}
\]

\[
\text{restriction theory}
\]

\[
\text{decoupling}
\]

\[
\text{time-frequency analysis}
\]

\[
\text{additive combinatorics}
\]

\[
\text{incidence geometry}
\]

\[
\text{ergodic averages}
\]

\[
\text{PDE dispersive}.
\]

Edges:

\[
f\mapsto \widehat f|_\Sigma
\]

\[
f\mapsto Tf
\]

\[
\text{curvature}\mapsto \text{decay/dispersion estimates}
\]

\[
\text{decoupling}\mapsto \text{Vinogradov-type counting}
\]

\[
\text{incidence bound}\mapsto \text{operator norm bound}.
\]

Scheme:

\[
\begin{gathered}
\text{curvature} \to \text{oscillation} \to \text{cancellation} \to \text{estimate} \\
\to \text{counting}.
\end{gathered}
\]



\subsection{D7. Network microlocal--symplectic--sheaf-theoretic--PDE}

Nodes:

\[
\text{distributions}
\]

\[
\text{wavefront sets}
\]

\[
\text{Lagrangiane coniche}
\]

\[
\text{Fourier integral operators}
\]

\[
\text{symplectic geometry}
\]

\[
\text{microlocal sheaves}
\]

\[
\text{Fukaya categories}
\]

\[
\text{PDE propagation}
\]

\[
\text{contact geometry}.
\]

Edges:

\[
u\mapsto WF(u)
\]

\[
P\mapsto H_{\sigma(P)}
\]

\[
\text{FIO}\mapsto \text{canonical relation}
\]

\[
\mathcal F\mapsto SS(\mathcal F)
\]

\[
\text{Lagrangian}\mapsto \text{sheaf quantization}.
\]

Idea:

\[
\begin{gathered}
\text{PDE singularities} \to \text{Legendrian geometry} \to \text{sheaf categories} \to \text{topological invariants}.
\end{gathered}
\]



\subsection{D8. Network spectral theory--scattering--zeta--traces}

Nodes:

\[
\text{operators autoaggiunti}
\]

\[
\text{non-self-adjoint operators}
\]

\[
\text{spectra discreti}
\]

\[
\text{spectra continui}
\]

\[
\text{resonances}
\]

\[
\text{scattering matrices}
\]

\[
\text{trace formulas}
\]

\[
\text{zeta functions}
\]

\[
\text{quantum chaos}
\]

\[
\text{inverse spectral geometry}.
\]

Edges:

\[
P\mapsto \operatorname{Spec}(P)
\]

\[
P\mapsto \zeta_P(s)
\]

\[
P\mapsto \operatorname{Tr}(e^{-tP})
\]

\[
\text{geodesic flow}\mapsto \text{trace formula}
\]

\[
V(x)\mapsto S_V(\lambda).
\]

Fibres:

\[
P\mapsto {\lambda_1,\ldots,\lambda_N}
\]

has very large fibres at low energy.



\subsection{D9. Network Riemannian geometry--Ricci flow--PDE geometric--topology}

Nodes:

\[
\text{metric riemanniane}
\]

\[
\text{curvature tensors}
\]

\[
\text{Ricci flow}
\]

\[
\text{mean curvature flow}
\]

\[
\text{Yamabe flow}
\]

\[
\text{geometric analysis}
\]

\[
\text{singularity models}
\]

\[
\text{topology of the varieties}
\]

\[
\text{limit metric spaces}.
\]

Edges:

\[
g\mapsto \operatorname{Rm}(g),\operatorname{Ric}(g),R(g)
\]

\[
g_0\mapsto g(t)
\]

\[
g(t)\mapsto \text{singularity model}
\]

\[
(M,g)\mapsto \text{volume growth, heat kernel, spectrum}
\]

\[
(M,g_i)\mapsto \text{Gromov--Hausdorff limit}.
\]

Scheme:

\[
\begin{gathered}
\text{geometry} \to \text{parabolic PDE} \to \text{monotonicity} \to \text{topology}.
\end{gathered}
\]



\subsection{D10. Network metric geometry--synthetic Ricci--optimal transport}

Nodes:

\[
\text{metric measure spaces}
\]

\[
\text{Wasserstein geometry}
\]

\[
\text{curvature-dimension conditions}
\]

\[
\text{optimal transport}
\]

\[
\text{gradient flows}
\]

\[
\text{entropy}
\]

\[
\text{heat flow}
\]

\[
\text{concentration}
\]

\[
\text{Markov chains}.
\]

Edges:

\[
(X,d,\mu)\mapsto \mathcal P_2(X)
\]

\[
\mu_0,\mu_1\mapsto \pi^\ast
\]

\[
\operatorname{Ent}_\mu\mapsto \text{gradient flow}
\]

\[
(X,d,\mu)\mapsto \text{CD}(K,N)\text{ profile}
\]

\[
\text{Markov chain}\mapsto \text{discrete transport geometry}.
\]

Scheme:

\[
\text{curvature}
\leftrightarrow
\text{convexity of the entropy}
\leftrightarrow
\text{heat contraction}
\leftrightarrow
\text{concentration}
\leftrightarrow
\text{mixing}.
\]



\subsection{D11. Network calculus of variations--\(\Gamma\)-convergence--materials--PDE}

Nodes:

\[
\text{functionals}
\]

\[
\text{minimizzatori}
\]

\[
\Gamma\text{-convergence}
\]

\[
\text{relaxation}
\]

\[
\text{quasiconvexity}
\]

\[
\text{Young measures}
\]

\[
\text{microstructures}
\]

\[
\text{elasticity}
\]

\[
\text{free boundary problems}
\]

\[
\text{optimal design}.
\]

Edges:

\[
F_\varepsilon\mapsto \Gamma\text{-}\lim F_\varepsilon
\]

\[
u_\varepsilon\mapsto \nu_x
\]

\[
\text{nonconvex energy}\mapsto \text{relaxed energy}
\]

\[
\text{material microstructure}\mapsto \text{effective tensor}
\]

\[
\text{shape}\mapsto \text{shape derivative}.
\]

Fundamental fibre:

\[
\text{microstructure}\mapsto \text{effective material}
\]

has very large fibres.



\subsection{D12. Network homogenization--multiscale--random media}

Nodes:

\[
\text{PDEs with oscillating coefficients}
\]

\[
\text{mezzi periodici}
\]

\[
\text{mezzi casuali}
\]

\[
H\text{-convergence}
\]

\[
G\text{-convergence}
\]

\[
\text{two-scale convergence}
\]

\[
\text{correctors}
\]

\[
\text{effective equations}
\]

\[
\text{percolation}
\]

\[
\text{random conductance models}.
\]

Edges:

\[
a(x/\varepsilon)\mapsto a_{\mathrm{hom}}
\]

\[
u_\varepsilon\mapsto u_0
\]

\[
u_\varepsilon\mapsto (u_0,u_1(x,y))
\]

\[
\omega\mapsto a_{\mathrm{hom}}(\omega)
\]

\[
\text{random medium}\mapsto \text{large-scale heat kernel}.
\]

Scheme:

\[
\text{micro}
\to
\text{meso}
\to
\text{macro}.
\]



\subsection{D13. Network fluids--turbulence--vorticity--statistics}

Nodes:

\[
\text{Euler}
\]

\[
\text{Navier--Stokes}
\]

\[
\text{vorticity}
\]

\[
\text{boundary layers}
\]

\[
\text{turbulenza}
\]

\[
\text{Reynolds averaging}
\]

\[
\text{large eddy simulation}
\]

\[
\text{Onsager regularity}
\]

\[
\text{statistical solutions}
\]

\[
\text{kinetic theory}.
\]

Edges:

\[
u\mapsto \omega=\nabla\times u
\]

\[
u^\nu\mapsto u^0
\]

\[
u\mapsto \overline u
\]

\[
u\mapsto E(k)
\]

\[
\text{particle system}\mapsto \text{kinetic equation}\mapsto \text{fluid equation}.
\]

Fibres:

\[
u\mapsto E(k)
\]

forgets phases and coherent vortices.

\[
u\mapsto \overline u
\]

forgets small scales.



\subsection{D14. Network kinetic theory--Boltzmann--entropy--hydrodynamic limits}

Nodes:

\[
\text{systems of particelle}
\]

\[
\text{Liouville equation}
\]

\[
\text{Boltzmann equation}
\]

\[
\text{Landau equation}
\]

\[
\text{Vlasov equation}
\]

\[
\text{Fokker--Planck}
\]

\[
\text{entropy methods}
\]

\[
\text{hydrodynamic limits}
\]

\[
\text{fluctuation theory}.
\]

Edges:

\[
{x_i,v_i}_{i=1}^N\mapsto f_N
\]

\[
f_N\mapsto f
\]

\[
f\mapsto \rho,u,T
\]

\[
\text{Boltzmann}\mapsto \text{Euler/Navier--Stokes}
\]

\[
f\mapsto H(f).
\]



\subsection{D15. Network continuous probability--Malliavin--SPDE--regularity}

Nodes:

\[
\text{Brownian motion}
\]

\[
\text{martingales}
\]

\[
\text{diffusions}
\]

\[
\text{Malliavin calculus}
\]

\[
\text{Dirichlet forms}
\]

\[
\text{SPDE}
\]

\[
\text{Gaussian fields}
\]

\[
\text{random distributions}
\]

\[
\text{stochastic quantization}.
\]

Edges:

\[
b,\sigma\mapsto X_t
\]

\[
X_t\mapsto L
\]

\[
L\mapsto \mathcal E
\]

\[
X\mapsto DX
\]

\[
\text{noise}\mapsto \text{random distribution}\mapsto \text{SPDE solution}.
\]



\subsection{D16. Network rough paths--regularity structures--renormalization}

Nodes:

\[
\text{rough paths}
\]

\[
\text{controlled paths}
\]

\[
\text{branched rough paths}
\]

\[
\text{regularity structures}
\]

\[
\text{paracontrolled distributions}
\]

\[
\text{renormalization group}
\]

\[
\text{Hopf algebras}
\]

\[
\text{SPDE singulari}
\]

\[
\text{stochastic analysis}.
\]

Edges:

\[
X\mapsto \mathbf X=(X,\mathbb X,\ldots)
\]

\[
\mathbf X\mapsto \text{solution map}
\]

\[
\text{SPDE}\mapsto \text{regularity structure}
\]

\[
\text{diagramma divergente}\mapsto \text{counterterm}
\]

\[
\text{trees}\mapsto \text{Hopf algebra}.
\]

Fundamental fibre:

\[
\mathbf X\mapsto X
\]

forgets iterated areas.

Idea:

\[
\text{segnale grezzo}
\not\Rightarrow
\text{dynamics ben posta}
\]

ma

\[
\text{segnale arricchito}
\Rightarrow
\text{dynamics continua}.
\]



\subsection{D17. Network ergodic theory--transfer operators--thermodynamic formalism}

Nodes:

\[
\text{systems dinamici misurabili}
\]

\[
\text{systems topologici}
\]

\[
\text{subshifts}
\]

\[
\text{operators of Koopman}
\]

\[
\text{operators of Perron--Frobenius}
\]

\[
\text{entropy}
\]

\[
\text{pressione}
\]

\[
\text{measures Gibbs}
\]

\[
\text{zeta dinamiche}
\]

\[
\text{large deviations}.
\]

Edges:

\[
T\mapsto U_T:f\mapsto f\circ T
\]

\[
T\mapsto \mathcal L_\phi
\]

\[
T\mapsto h_\mu(T),h_{\mathrm{top}}(T)
\]

\[
T\mapsto \zeta_T(z)
\]

\[
\text{partition}\mapsto \text{symbolic coding}.
\]



\subsection{D18. Network chaos--KAM--small divisors--localization}

Nodes:

\[
\text{systems hamiltoniani}
\]

\[
\text{tori invariants}
\]

\[
\text{KAM theory}
\]

\[
\text{small divisors}
\]

\[
\text{Diophantine approximation}
\]

\[
\text{quasi-periodic Schrödinger operators}
\]

\[
\text{Anderson localization}
\]

\[
\text{Aubry--Mather theory}
\]

\[
\text{renormalization}.
\]

Edges:

\[
H(I,\theta)\mapsto \omega(I)
\]

\[
\omega\mapsto \text{Diophantine exponent}
\]

\[
\text{perturbation}\mapsto \text{surviving tori}
\]

\[
\text{quasi-periodic potential}\mapsto \text{spectrum}
\]

\[
\text{cocycle}\mapsto \text{Lyapunov exponent}.
\]

Idea:

\[
\text{arithmetic fine}
\to
\text{stability dynamics}
\to
\text{spectral analysis}.
\]



\subsection{D19. Network integrable systems--Riemann--Hilbert--isomonodromy}

Nodes:

\[
\text{equazioni integrabili}
\]

\[
\text{Lax pairs}
\]

\[
\text{spectral curves}
\]

\[
\text{inverse scattering}
\]

\[
\text{Riemann--Hilbert problems}
\]

\[
\text{isomonodromic deformations}
\]

\[
\text{Painlevé equations}
\]

\[
\text{tau functions}
\]

\[
\text{Frobenius manifolds}
\]

\[
\text{random matrix theory}.
\]

Edges:

\[
\partial_t L=[P,L]\mapsto \operatorname{Spec}(L)
\]

\[
u(x)\mapsto \text{scattering data}
\]

\[
\text{scattering data}\mapsto \text{Riemann--Hilbert problem}
\]

\[
\text{monodromy data}\mapsto \text{isomonodromic flow}
\]

\[
\text{RHP}\mapsto \tau.
\]



\subsection{D20. Network complex analysis--SCV--pluripotential--PDE}

Nodes:

\[
\text{holomorphic functions}
\]

\[
\text{pseudoconvex domains}
\]

\[
\bar\partial\text{-problem}
\]

\[
\text{Bergman spaces}
\]

\[
\text{Hardy spaces}
\]

\[
\text{plurisubharmonic functions}
\]

\[
\text{Monge--Ampère complex}
\]

\[
\text{Stein manifolds}
\]

\[
\text{complex dynamics}
\]

\[
\text{Kähler geometry}.
\]

Edges:

\[
\Omega\mapsto K_\Omega(z,w)
\]

\[
\Omega\mapsto \bar\partial\text{-Neumann operator}
\]

\[
\varphi\mapsto (dd^c\varphi)^n
\]

\[
X\mapsto H^{p,q}_{\bar\partial}(X)
\]

\[
f:\mathbb C^n\to\mathbb C^n\mapsto \text{Green current}.
\]



\subsection{D21. Network potential--capacity--electrical networks--probability}

Nodes:

\[
\text{potential classico}
\]

\[
\text{capacity}
\]

\[
\text{Dirichlet energy}
\]

\[
\text{electrical networks}
\]

\[
\text{random walks}
\]

\[
\text{Brownian motion}
\]

\[
\text{harmonic measure}
\]

\[
\text{Laplacian growth}
\]

\[
\text{potential theory on graphs}.
\]

Edges:

\[
D\mapsto \operatorname{cap}(D)
\]

\[
f\mapsto \int|\nabla f|^2
\]

\[
G\mapsto \text{electrical network}
\]

\[
\text{network}\mapsto \text{effective resistance metric}
\]

\[
\text{domain}\mapsto \omega^x_{\partial D}.
\]



\subsection{D22. Network non-self-adjoint operators--pseudospectra--numerical stability}

Nodes:

\[
\text{operators non normali}
\]

\[
\text{pseudospectra}
\]

\[
\text{semigroups}
\]

\[
\text{linear stability}
\]

\[
\text{transient growth}
\]

\[
\text{fluid stability}
\]

\[
\text{numerical analysis}
\]

\[
\text{control}.
\]

Edges:

\[
A\mapsto \operatorname{Spec}(A)
\]

\[
A\mapsto |(z-A)^{-1}|
\]

\[
A\mapsto \operatorname{Pseudospec}_\varepsilon(A)
\]

\[
A\mapsto e^{tA}
\]

\[
A\mapsto \sup_t|e^{tA}|.
\]

Principle:

\[
\text{same spectrum}
\not\Rightarrow
\text{same dynamics}.
\]



\subsection{D23. Network approximation--wavelet--compressed sensing--numerical learning}

Nodes:

\[
\text{approximation theory}
\]

\[
\text{polynomial approximation}
\]

\[
\text{splines}
\]

\[
\text{wavelets}
\]

\[
\text{sparse approximation}
\]

\[
\text{compressed sensing}
\]

\[
\text{reduced basis methods}
\]

\[
\text{neural operators}
\]

\[
\text{numerical PDE}
\]

\[
\text{information-based complexity}.
\]

Edges:

\[
f\mapsto \text{best }n\text{-term approximation}
\]

\[
f\mapsto {\langle f,\psi_\lambda\rangle}
\]

\[
u_\mu\mapsto \text{reduced basis space}
\]

\[
Au=f\mapsto A_hu_h=f_h
\]

\[
\text{operator}\mapsto \text{learned operator surrogate}.
\]

Scheme:

\[
\begin{gathered}
\text{regularity} \to \text{compressibility} \to \text{sampling} \to \text{algorithms} \\
\to \text{error verification data}.
\end{gathered}
\]



\subsection{D24. Network numerical analysis--FEM--mesh--verification data}

Nodes:

\[
\text{finite elements}
\]

\[
\text{finite volumes}
\]

\[
\text{spectral methods}
\]

\[
\text{mesh refinement}
\]

\[
\text{adaptive algorithms}
\]

\[
\text{a posteriori estimates}
\]

\[
\text{verified numerics}
\]

\[
\text{interval arithmetic}
\]

\[
\text{computer-assisted proofs}.
\]

Edges:

\[
\text{PDE}\mapsto \text{weak formulation}
\]

\[
\text{weak formulation}\mapsto \text{discrete variational problem}
\]

\[
u_h\mapsto \eta_h
\]

\[
\eta_h\mapsto \text{mesh refinement}
\]

\[
\text{floating computation}\mapsto \text{interval verification data}.
\]

Scheme:

\[
\begin{gathered}
\text{analysis} \to \text{algorithm} \to \text{verification datum} \to \text{formalization}.
\end{gathered}
\]



\subsection{D25. Network control--observability--identification--optimization}

Nodes:

\[
\text{controlled dynamical systems}
\]

\[
\text{ODE}
\]

\[
\text{PDE controllate}
\]

\[
\text{Hamilton--Jacobi--Bellman}
\]

\[
\text{Pontryagin maximum principle}
\]

\[
\text{reachability}
\]

\[
\text{observability}
\]

\[
\text{system identification}
\]

\[
\text{mean field games}
\]

\[
\text{reinforcement learning}.
\]

Edges:

\[
\dot x=f(x,u)\mapsto \mathcal R_T(x_0)
\]

\[
\text{control problem}\mapsto \text{value function}
\]

\[
\text{value function}\mapsto \text{HJB equation}
\]

\[
\text{trajectory data}\mapsto \hat f
\]

\[
\text{multi-agent system}\mapsto \text{mean field limit}.
\]

Fibre:

\[
f\mapsto \text{input-output behavior}
\]

may have very large fibres.



\subsection{D26. Network functional analysis--Banach--operator ideals--asymptotic geometry}

Nodes:

\[
\text{Banach spaces}
\]

\[
\text{Hilbert spaces}
\]

\[
\text{operator ideals}
\]

\[
\text{type/cotype}
\]

\[
\text{local theory}
\]

\[
\text{basis constants}
\]

\[
\text{compact operators}
\]

\[
\text{nuclear operators}
\]

\[
\text{tensor norms}
\]

\[
\text{nonlinear embeddings}.
\]

Edges:

\[
X\mapsto {E\subset X:\dim E\le n}
\]

\[
T:X\to Y\mapsto (s_n(T))
\]

\[
X\mapsto \text{type/cotype profile}
\]

\[
X,Y\mapsto X\widehat\otimes_\pi Y,\ X\widehat\otimes_\varepsilon Y
\]

\[
X\mapsto \operatorname{Lip}_0(X).
\]



\subsection{D27. Network interpolation--scales of spaces--adaptive regularity}

Nodes:

\[
L^p
\]

\[
\text{Sobolev}
\]

\[
\text{Besov}
\]

\[
\text{Triebel--Lizorkin}
\]

\[
\text{Hardy}
\]

\[
\text{BMO}
\]

\[
\text{Lorentz}
\]

\[
\text{Orlicz}
\]

\[
\text{interpolation spaces}
\]

\[
\text{regularity estimates}.
\]

Edges:

\[
(X_0,X_1)\mapsto [X_0,X_1]_\theta
\]

\[
(X_0,X_1)\mapsto (X_0,X_1)_{\theta,q}
\]

\[
f\mapsto \text{regularity profile }s\mapsto |f|_{X_s}
\]

\[
T:X_0\to Y_0,\ T:X_1\to Y_1
\mapsto
T:X_\theta\to Y_\theta.
\]



\subsection{D28. Network analytic nonlocality--fractional operators--jump processes}

Nodes:

\[
\text{fractional Laplacians}
\]

\[
\text{nonlocal PDE}
\]

\[
\text{integro-differential operators}
\]

\[
\text{stable processes}
\]

\[
\text{nonlocal minimal surfaces}
\]

\[
\text{peridynamics}
\]

\[
\text{anomalous diffusion}.
\]

Edges:

\[
(-\Delta)^s\mapsto \text{stable Lévy process}
\]

\[
u\mapsto \int\frac{u(x)-u(y)}{|x-y|^{n+2s}},dy
\]

\[
s\to 1
\]

\[
s\to 0
\]

\[
K(x,y)\mapsto \text{jump process}.
\]



\subsection{D29. Network memory--delay equations--history operators--systems}

Nodes:

\[
\text{delay differential equations}
\]

\[
\text{Volterra equations}
\]

\[
\text{fractional time derivatives}
\]

\[
\text{memory kernels}
\]

\[
\text{viscoelasticity}
\]

\[
\text{control with memory}
\]

\[
\text{non-Markovian processes}
\]

\[
\text{state augmentation}.
\]

Edges:

\[
K(t)\mapsto \int_0^t K(t-s)u(s),ds
\]

\[
\text{memory system}\mapsto \text{augmented Markovian system}
\]

\[
K\mapsto \widehat K(\lambda)
\]

\[
\text{fractional derivative}\mapsto \text{subordinated semigroup}.
\]

Idea:

\[
\text{the presentation of a system may include the amount of past needed}.
\]



\subsection{D30. Network algebraic analysis--D-modules--PDE--perverse sheaves}

Nodes:

\[
D\text{-modules}
\]

\[
\text{linear PDE systems}
\]

\[
\text{holonomic modules}
\]

\[
\text{characteristic varieties}
\]

\[
\text{Riemann--Hilbert correspondence}
\]

\[
\text{perverse sheaves}
\]

\[
\text{monodromy}
\]

\[
\text{Fourier--Laplace transform}
\]

\[
\text{irregular singularities}
\]

\[
\text{Stokes data}.
\]

Edges:

\[
\text{PDE system}\mapsto D/D\langle P_i\rangle
\]

\[
M\mapsto \operatorname{Char}(M)
\]

\[
M_{\mathrm{hol}}\mapsto \operatorname{Sol}(M)
\]

\[
\operatorname{Sol}(M)\mapsto \text{perverse sheaf}
\]

\[
\text{irregular connection}\mapsto \text{Stokes data}.
\]

Scheme:

\[
\text{linear PDEs}
\leftrightarrow
\text{noncommutative algebra}
\leftrightarrow
\text{constructible topology}.
\]



\subsection{D31. Network transseries--asymptotics--resurgence--ODE/PDE}

Nodes:

\[
\text{asymptotic expansions}
\]

\[
\text{transseries}
\]

\[
\text{resurgence}
\]

\[
\text{Borel summation}
\]

\[
\text{Stokes phenomena}
\]

\[
\text{singular ODE}
\]

\[
\text{WKB analysis}
\]

\[
\text{quantum curves}
\]

\[
\text{enumerative geometry}.
\]

Edges:

\[
f(\varepsilon)\mapsto \sum a_n\varepsilon^n
\]

\[
\sum a_n\varepsilon^n\mapsto \mathcal Bf
\]

\[
\text{Borel singularities}\mapsto \text{Stokes constants}
\]

\[
\text{semiclassical ODE}\mapsto \text{WKB transseries}
\]

\[
\text{enumerative generating function}\mapsto \text{resurgent structure}.
\]

Idea:

\[
\begin{gathered}
\text{funzione} \to \text{serie divergente} \to \text{Borel singularities} \to \text{Stokes data} \\
\to \text{monodromy}.
\end{gathered}
\]



\subsection{D32. Network analytic QFT--renormalization--operator algebras--probability}

Nodes:

\[
\text{Euclidean QFT}
\]

\[
\text{constructive field theory}
\]

\[
\text{renormalization group}
\]

\[
\text{Gaussian free fields}
\]

\[
\text{CFT}
\]

\[
\text{OPE algebras}
\]

\[
\text{stochastic quantization}
\]

\[
\text{AQFT}
\]

\[
\text{factorization algebras}
\]

\[
\text{perturbative QFT}.
\]

Edges:

\[
S\mapsto e^{-S}\mathcal D\phi
\]

\[
\text{field}\mapsto \text{correlation functions}
\]

\[
\text{correlators}\mapsto \text{OPE coefficients}
\]

\[
\text{cutoff theory}\mapsto \text{renormalized theory}
\]

\[
\text{local net}\mapsto \text{operator algebra}.
\]

Idea:

QFT has many presentations: Lagrangian, Hamiltonian, operator-algebraic, correlational, and categorical.



\subsection{D33. Network causality--Lorentzian geometry--hyperbolic PDE--relativity}

Nodes:

\[
\text{Lorentzian manifolds}
\]

\[
\text{causal structures}
\]

\[
\text{Einstein equations}
\]

\[
\text{wave equations}
\]

\[
\text{null geodesics}
\]

\[
\text{Penrose diagrams}
\]

\[
\text{conformal geometry}
\]

\[
\text{inverse problems}
\]

\[
\text{black hole scattering}.
\]

Edges:

\[
g\mapsto \text{light cones}
\]

\[
g\mapsto \square_g
\]

\[
\square_g\mapsto \text{propagator}
\]

\[
\text{null geodesic flow}\mapsto \text{lens relation}
\]

\[
\text{Einstein equations}\mapsto \text{constraint equations}.
\]



\subsection{D34. Network condensed mathematics--derived functional analysis}

Nodes:

\[
\text{topological vector spaces}
\]

\[
\text{locally convex spaces}
\]

\[
\text{Fréchet spaces}
\]

\[
\text{nuclear spaces}
\]

\[
\text{Schwartz distributions}
\]

\[
\text{condensed mathematics}
\]

\[
\text{solid modules}
\]

\[
\text{derived functional analysis}
\]

\[
\text{analytic geometry over rings}.
\]

Edges:

\[
V\mapsto \underline V
\]

\[
\text{functional analytic problem}\mapsto \text{abelian categorical problem}
\]

\[
\text{distribution space}\mapsto \text{sheaf/condensed module}
\]

\[
\text{topological tensor product}\mapsto \text{derived tensor}.
\]

Idea:

\[
\text{absorbing functional analysis into a stable categorical language}.
\]



\section{E. Probability, Statistics, Information, Learning, and Applications}



\subsection{E1. Network probability--statistics--information--algebraic geometry}

Nodes:

\[
\text{distributions discrete}
\]

\[
\text{measures continue}
\]

\[
\text{moments}
\]

\[
\text{exponential families}
\]

\[
\text{graphical models}
\]

\[
\text{algebraic statistics}
\]

\[
\text{toric varieties}
\]

\[
\text{information geometry}
\]

\[
\text{optimal transport}
\]

\[
\text{learning theory}
\]

\[
\text{coding theory}.
\]

Edges:

\[
p\mapsto (m_\alpha(p))_{|\alpha|\le d}
\]

\[
\text{graphical model}\mapsto \text{conditional independence ideal}
\]

\[
\text{Bayesian network}\mapsto \text{semialgebraic model}
\]

\[
\mu,\nu\mapsto W_p(\mu,\nu)
\]

\[
\text{code}\mapsto \text{probability channel}.
\]

Central fibre:

\[
\theta\mapsto p_\theta
\]

measure identifiability and non-identifiability.



\subsection{E2. Network information--entropy--compression--analysis}

Nodes:

\[
\text{Shannon entropy}
\]

\[
\text{Rényi entropy}
\]

\[
\text{Fisher information}
\]

\[
\text{log-Sobolev inequalities}
\]

\[
\text{transport inequalities}
\]

\[
\text{compression}
\]

\[
\text{statistical estimation}
\]

\[
\text{large deviations}
\]

\[
\text{thermodynamic formalism}
\]

\[
\text{quantum information}.
\]

Edges:

\[
\mu\mapsto H(\mu)
\]

\[
\mu\mapsto I(\mu)
\]

\[
\text{Markov semigroup}\mapsto \text{entropy dissipation}
\]

\[
\text{source}\mapsto \text{code length profile}
\]

\[
\rho\mapsto S(\rho).
\]

Scheme:

\[
\text{entropy}
\leftrightarrow
\text{heat}
\leftrightarrow
\text{optimal transport}
\leftrightarrow
\text{concentration}
\leftrightarrow
\text{compression}
\leftrightarrow
\text{learning}.
\]



\subsection{E3. Network large deviations--variational principles--limit PDE}

Nodes:

\[
\text{probability measures}
\]

\[
\text{large deviation principles}
\]

\[
\text{rate functions}
\]

\[
\text{Hamilton--Jacobi equations}
\]

\[
\text{statistical mechanics}
\]

\[
\text{mean field limits}
\]

\[
\text{optimal control}
\]

\[
\text{rare events}.
\]

Edges:

\[
X_\varepsilon\mapsto I
\]

\[
I\mapsto \text{variational problem}
\]

\[
\text{Markov process}\mapsto \text{Hamiltonian}
\]

\[
\text{rare event}\mapsto \text{optimal path}.
\]



\subsection{E4. Network statistical mechanics--Gibbs measures--phase transitions--complexity}

Nodes:

\[
\text{spin systems}
\]

\[
\text{Gibbs measures}
\]

\[
\text{partition functions}
\]

\[
\text{cluster expansions}
\]

\[
\text{phase transitions}
\]

\[
\text{random fields}
\]

\[
\text{percolation}
\]

\[
\text{Ising/Potts models}
\]

\[
\text{statistical algorithms}
\]

\[
\text{complexity of counting}.
\]

Edges:

\[
H(\sigma)\mapsto Z(\beta)
\]

\[
H\mapsto \mu_\beta
\]

\[
\mu_\beta\mapsto \text{correlation functions}
\]

\[
\text{spin model}\mapsto \text{graphical expansion}
\]

\[
\text{Gibbs sampler}\mapsto \text{Markov chain}.
\]



\subsection{E5. Network mathematical machine learning--kernels--operators--geometry}

Nodes:

\[
\text{kernel methods}
\]

\[
\text{RKHS}
\]

\[
\text{Gaussian processes}
\]

\[
\text{neural networks}
\]

\[
\text{neural tangent kernels}
\]

\[
\text{mean field limits}
\]

\[
\text{optimal transport}
\]

\[
\text{statistical learning theory}
\]

\[
\text{PAC-Bayes}
\]

\[
\text{information bottleneck}.
\]

Edges:

\[
k(x,y)\mapsto \mathcal H_k
\]

\[
\text{network}\mapsto \text{function class}
\]

\[
\text{wide network}\mapsto \text{kernel/mean-field limit}
\]

\[
\text{data distribution}\mapsto \text{risk functional}
\]

\[
\text{training dynamics}\mapsto \text{gradient flow}.
\]

Central fibre:

\[
\text{parameters}\mapsto \text{function}
\]

has very large fibres.



\subsection{E6. Network scientific machine learning--PDE discovery--operator learning}

Nodes:

\[
\text{PDE data}
\]

\[
\text{symbolic regression}
\]

\[
\text{operator learning}
\]

\[
\text{neural operators}
\]

\[
\text{physics-informed neural networks}
\]

\[
\text{reduced order models}
\]

\[
\text{Bayesian calibration}
\]

\[
\text{uncertainty quantification}
\]

\[
\text{digital twins}.
\]

Edges:

\[
\text{data }u(x,t)\mapsto \hat P
\]

\[
P\mapsto \mathcal S_P
\]

\[
\mathcal S_P\mapsto \widehat{\mathcal S}
\]

\[
\text{surrogate}\mapsto \text{certified error bound}
\]

\[
\text{physical constraints}\mapsto \text{loss terms}.
\]

Fibre:

\[
\text{observed data}
\mapsto
\text{insieme of the PDE compatibili}.
\]



\subsection{E7. Network inverse problems--tomography--machine learning--invariants}

Nodes:

\[
\text{signals}
\]

\[
\text{immagini}
\]

\[
\text{varieties}
\]

\[
\text{graphs}
\]

\[
\text{groups}
\]

\[
\text{Radon transforms}
\]

\[
\text{scattering transforms}
\]

\[
\text{graph neural networks}
\]

\[
\text{kernel methods}
\]

\[
\text{statistical queries}
\]

\[
\text{invariant theory}.
\]

Edges:

\[
X\mapsto \mathcal O_{\le r}(X)
\]

\[
G\curvearrowright X\mapsto k[X]^G_{\le d}
\]

\[
\text{graph}\mapsto \text{message-passing features}
\]

\[
\text{signal}\mapsto \text{wavelet/scattering coefficients}
\]

\[
\text{measure}\mapsto \text{statistical query oracle}.
\]

Question:

\[
\text{what can a bounded class of observers see?}
\]



\subsection{E8. Network universal inverse problems}

Nodes:

\[
\text{object nascosto}
\]

\[
\text{forward map}
\]

\[
\text{observed data}
\]

\[
\text{regularization}
\]

\[
\text{stability estimates}
\]

\[
\text{Bayesian inverse problems}
\]

\[
\text{optimal experimental design}
\]

\[
\text{compressed sensing}
\]

\[
\text{tomography}
\]

\[
\text{learning-based inversion}.
\]

Arco fondamentale:

\[
x\mapsto F(x).
\]

Examples:

\[
V\mapsto \text{scattering data}
\]

\[
\text{domain}\mapsto \text{boundary measurements}
\]

\[
\text{immagine}\mapsto \text{Radon transform}
\]

\[
\text{metric}\mapsto \text{lens data}
\]

\[
\text{coefficiente PDE}\mapsto \text{Dirichlet-to-Neumann map}.
\]

Principle:

\[
\text{how much does it cost to distinguish two hidden objects through accessible data?}
\]



\subsection{E9. Network mathematical economics--equilibria--convexity--dynamics}

Nodes:

\[
\text{fixed point theory}
\]

\[
\text{game theory}
\]

\[
\text{Nash equilibria}
\]

\[
\text{market equilibria}
\]

\[
\text{mechanism design}
\]

\[
\text{optimal transport economics}
\]

\[
\text{mean field games}
\]

\[
\text{variational inequalities}
\]

\[
\text{convex analysis}
\]

\[
\text{algorithmic game theory}.
\]

Edges:

\[
\text{game}\mapsto \text{best response correspondence}
\]

\[
\text{best response}\mapsto \text{fixed point problem}
\]

\[
\text{market}\mapsto \text{convex program}
\]

\[
\text{many-agent game}\mapsto \text{mean field game PDE}
\]

\[
\text{mechanism}\mapsto \text{incentive constraints}.
\]



\subsection{E10. Network mathematical biology--reaction networks--dynamics--topology}

Nodes:

\[
\text{chemical reaction networks}
\]

\[
\text{ODE systems}
\]

\[
\text{stochastic reaction networks}
\]

\[
\text{mass-action kinetics}
\]

\[
\text{Petri nets}
\]

\[
\text{persistence}
\]

\[
\text{bifurcation theory}
\]

\[
\text{network motifs}
\]

\[
\text{phylogenetics}
\]

\[
\text{algebraic statistics}.
\]

Edges:

\[
\text{reaction network}\mapsto \text{ODE}
\]

\[
\text{reaction network}\mapsto \text{Petri net}
\]

\[
\text{ODE}\mapsto \text{steady-state ideal}
\]

\[
\text{stochastic network}\mapsto \text{master equation}
\]

\[
\text{phylogenetic tree}\mapsto \text{statistical moof the variety}.
\]



\section{F. Arithmetic, Automata, Languages, Discrete Mathematics, and Computability}



\subsection{F1. Network computability--tilings--subshift--groups--cellular automata}

Nodes:

\[
\text{macchine of Turing}
\]

\[
\text{tag systems}
\]

\[
\text{Wang tilings}
\]

\[
\text{subshifts of finite type}
\]

\[
\text{cellular automata}
\]

\[
\text{semifinitely presented groups}
\]

\[
\text{finitely presented groups}
\]

\[
\text{rewriting systems}
\]

\[
\text{moof the checking}
\]

\[
\text{reachability}.
\]

Edges:

\[
\text{Turing machine}\mapsto \text{tiling system}
\]

\[
\text{tiling}\mapsto \text{subshift}
\]

\[
\text{subshift}\mapsto \text{cellular automaton}
\]

\[
\text{rewriting}\mapsto \text{semigroup presentation}
\]

\[
\text{semigroup}\mapsto \text{group embedding}.
\]

Effetto:

\[
\text{undecidability}
\to
\text{noncomputable quantitative profiles}.
\]



\subsection{F2. Network arithmetic--Galois--monodromy--automata--languages}

Nodes:

\[
\text{campi of numeri}
\]

\[
\text{arithmetic varieties}
\]

\[
\text{representations of Galois}
\]

\[
\ell\text{-adic cohomology}
\]

\[
\text{equazioni differentials}
\]

\[
\text{monodromy}
\]

\[
\text{dessins of enfants}
\]

\[
\text{ricorrenze lineari}
\]

\[
\text{sequences automatiche}
\]

\[
\text{formal languages}
\]

\[
\text{compression}.
\]

Edges:

\[
X/\mathbb Q\mapsto \rho_{X,\ell}
\]

\[
X\mapsto \#X(\mathbb F_p)
\]

\[
\text{covering}\mapsto \text{monodromy permutation group}
\]

\[
\text{linear recurrence}\mapsto \text{companion matrix}
\]

\[
\alpha\mapsto \text{digit word}_b(\alpha)
\]

\[
\text{sequence}\mapsto \text{automaton/minimal transducer}.
\]

Questions:

\[
\text{how much arithmetic is visible from a digit prefix?}
\]

\[
\text{how much geometry is visible from counts modulo small primes?}
\]



\subsection{F3. Network automorphic forms--L-functions--spectral PDE--arithmetic}

Nodes:

\[
\text{modular forms}
\]

\[
\text{automorphic forms}
\]

\[
\text{representations automorfe}
\]

\[
L\text{-functions}
\]

\[
\text{trace formulas}
\]

\[
\text{Shimura varieties}
\]

\[
\text{Galois representations}
\]

\[
\text{spectral theory on locally symmetric spaces}
\]

\[
\text{periods}.
\]

Edges:

\[
f\mapsto L(f,s)
\]

\[
f\mapsto (a_n(f))_{n\le N}
\]

\[
\pi\mapsto \rho_{\pi,\ell}
\]

\[
\Gamma\backslash G/K\mapsto \operatorname{Spec}(\Delta)
\]

\[
\text{trace formula}\mapsto \text{comparison of spectra}.
\]



\subsection{F4. Network formal languages--automata--spectral analysis--symbolic dynamics}

Nodes:

\[
\text{regular languages}
\]

\[
\text{context-free languages}
\]

\[
\text{weighted automata}
\]

\[
\text{formal power series}
\]

\[
\text{subshifts}
\]

\[
\text{transfer matrices}
\]

\[
\text{thermodynamic formalism}
\]

\[
\text{zeta functions}
\]

\[
\text{compression}
\]

\[
\text{random walks on automata}.
\]

Edges:

\[
\mathcal A\mapsto L(\mathcal A)
\]

\[
\mathcal A\mapsto A_{\mathcal A}
\]

\[
A_{\mathcal A}\mapsto \operatorname{Spec}(A_{\mathcal A})
\]

\[
L\mapsto \sum_n a_nz^n
\]

\[
\text{subshift}\mapsto \zeta(z).
\]

Scheme:

\[
\begin{gathered}
\text{languages} \to \text{matrices} \to \text{spectra} \to \text{dynamics} \\
\to \text{analysis}.
\end{gathered}
\]



\subsection{F5. Network finite observable profiles: quotients, truncations, reductions, completions}

Nodes:

\[
\text{objects infiniti}
\]

\[
\text{finite quotients}
\]

\[
\text{jets}
\]

\[
\text{truncations}
\]

\[
\text{reductions mod }p
\]

\[
\text{completions profiniti}
\]

\[
p\text{-adic completions}
\]

\[
\text{localizations}
\]

\[
\text{finite-dimensional approximations}.
\]

Edges:

\[
X\mapsto X_{\le r}
\]

\[
G\mapsto G/N
\]

\[
X/\mathbb Z\mapsto X_{\mathbb F_p}
\]

\[
R\mapsto R/\mathfrak m^k
\]

\[
A\mapsto A_n
\]

\[
X\mapsto \widehat X.
\]

Idea:

\[
\begin{gathered}
\text{every theory has finite observable profiles;}\\
\text{complexity is the profile by which they reveal the object.}
\end{gathered}
\]



\subsection{F6. Network codes--lattices--arithmetic--mathematical cryptography}

Nodes:

\[
\text{lattices}
\]

\[
\text{linear codes}
\]

\[
\text{codes algebro-geometrici}
\]

\[
\text{geometry of numbers}
\]

\[
\text{isogeny graphs}
\]

\[
\text{modular forms}
\]

\[
\text{finite fields}
\]

\[
\text{group actions}
\]

\[
\text{hard homogeneous spaces}
\]

\[
\text{proof systems}.
\]

Edges:

\[
C\mapsto \Lambda_C
\]

\[
X/\mathbb F_q\mapsto C_X
\]

\[
\text{isogeny graph}\mapsto \text{expander/random walk}
\]

\[
\text{lattice}\mapsto \text{theta series}
\]

\[
\text{constraint system}\mapsto \text{interactive proof object}.
\]

Observables:

\[
[n,k,d]
\]

\[
\text{theta series}
\]

\[
\text{spectrum}
\]

\[
\text{weight enumerator}.
\]



\subsection{F7. Network symmetry--species--generating functions--representations}

Nodes:

\[
\text{combinatorial species}
\]

\[
\text{operadi}
\]

\[
\text{symmetric functions}
\]

\[
S_n\text{-representations}
\]

\[
\text{Hilbert schemes}
\]

\[
\text{moduli spaces}
\]

\[
\text{generating functions}
\]

\[
D\text{-finite functions}
\]

\[
\text{recurrences}
\]

\[
\text{topological recursion}.
\]

Edges:

\[
\mathcal S\mapsto F_{\mathcal S}(x)
\]

\[
\mathcal S\mapsto Z_{\mathcal S}
\]

\[
X_n\mapsto H^\ast(X_n)
\]

\[
\text{sequence}\mapsto \text{linear recurrence/differential equation}
\]

\[
\text{moduli problem}\mapsto \text{generating series}.
\]

Question:

\[
\text{how much of a structure is visible from its generating function?}
\]



\subsection{F8. Network geometric group theory--analysis--spectra--probability}

Nodes:

\[
\text{groups finitamente generati}
\]

\[
\text{Cayley graphs}
\]

\[
\text{random walks}
\]

\[
\text{growth functions}
\]

\[
\text{isoperimetric profiles}
\]

\[
\text{property T}
\]

\[
\text{coarse embeddings}
\]

\[
\text{boundaries}
\]

\[
\text{operator algebras}
\]

\[
L^2\text{-cohomology}.
\]

Edges:

\[
(G,S)\mapsto \operatorname{Cay}(G,S)
\]

\[
G\mapsto \gamma_G(n)
\]

\[
G\mapsto \text{return probabilities}
\]

\[
G\mapsto C^\ast_r(G)
\]

\[
G\mapsto \beta_i^{(2)}(G).
\]



\subsection{F9. Network discrete-continuous: graphs, meshes, manifolds, operators}

Nodes:

\[
\text{graphs}
\]

\[
\text{simplicial complexes}
\]

\[
\text{meshes}
\]

\[
\text{manifolds}
\]

\[
\text{discrete Laplacians}
\]

\[
\text{finite element operators}
\]

\[
\text{graph neural operators}
\]

\[
\text{spectral convergence}
\]

\[
\text{Gromov--Hausdorff convergence}.
\]

Edges:

\[
M\mapsto \mathcal T_h(M)
\]

\[
\mathcal T_h\mapsto L_h
\]

\[
L_h\mapsto \operatorname{Spec}(L_h)
\]

\[
G_n\mapsto M
\]

\[
\text{graph signal}\mapsto \text{continuum function}.
\]



\subsection{F10. Network graph limits--continuous analysis--extremal combinatorics}

Nodes:

\[
\text{finite graphs}
\]

\[
\text{graphons}
\]

\[
\text{hypergraphons}
\]

\[
\text{flag algebras}
\]

\[
\text{cut norm}
\]

\[
\text{exchangeable arrays}
\]

\[
\text{dense graph limits}
\]

\[
\text{sparse graph limits}
\]

\[
\text{property testing}.
\]

Edges:

\[
G_n\mapsto W
\]

\[
G\mapsto (t(F,G))_F
\]

\[
W\mapsto T_W
\]

\[
W\mapsto \operatorname{Spec}(T_W)
\]

\[
\text{exchangeable array}\mapsto \text{graphon representation}.
\]



\section{G. Quantum and Noncommutative Networks}



\subsection{G1. Network quantum--noncommutative--operatorial--tensor networks}

Nodes:

\[
\text{quantum states}
\]

\[
\text{quantum channels}
\]

\[
\text{local Hamiltonians}
\]

\[
\text{tensor networks}
\]

\[
\text{stabilizer formalism}
\]

\[
C^\ast\text{-algebras}
\]

\[
\text{von Neumann algebras}
\]

\[
\text{nonlocal games}
\]

\[
\text{semidefinite programming}
\]

\[
\text{noncommutative polynomial optimization}
\]

\[
\text{modular tensor categories}.
\]

Edges:

\[
\text{nonlocal game}\mapsto \text{operator relations}
\]

\[
H=\sum H_i\mapsto \text{quantum CSP}
\]

\[
\psi\mapsto \text{tensor network ansatz}
\]

\[
\text{operator inequalities}\mapsto \text{SDP hierarchy}
\]

\[
\text{TQFT}\mapsto \text{quantum circuit representation}.
\]



\subsection{G2. Network quantum many-body--limit PDE--tensor networks--operators}

Nodes:

\[
\text{many-body quantum systems}
\]

\[
\text{local Hamiltonians}
\]

\[
\text{mean field limits}
\]

\[
\text{Hartree equation}
\]

\[
\text{Gross--Pitaevskii}
\]

\[
\text{density matrices}
\]

\[
\text{BBGKY hierarchy}
\]

\[
\text{tensor networks}
\]

\[
\text{entanglement entropy}
\]

\[
\text{operator algebras}.
\]

Edges:

\[
\Psi_N\mapsto \gamma_N^{(k)}
\]

\[
\gamma_N^{(k)}\mapsto \gamma^{\otimes k}
\]

\[
\Psi_N\mapsto \text{Hartree/NLS limit}
\]

\[
\Psi\mapsto \text{MPS/PEPS approximation}
\]

\[
H\mapsto \text{ground state}\mapsto \text{entanglement profile}.
\]

Fibre:

\[
\Psi\mapsto {\gamma^{(k)}}_{k\le r}
\]

forgets higher correlations.



\subsection{G3. Network operator algebras--index--K-theory--coarse geometry}

Nodes:

\[
C^\ast\text{-algebras}
\]

\[
\text{von Neumann algebras}
\]

\[
\text{Fredholm modules}
\]

\[
\text{spectral triples}
\]

\[
K\text{-theory}
\]

\[
K\text{-homology}
\]

\[
\text{index theory}
\]

\[
\text{coarse geometry}
\]

\[
\text{group }C^\ast\text{-algebras}
\]

\[
\text{noncommutative geometry}.
\]

Edges:

\[
X\mapsto C(X)
\]

\[
\Gamma\mapsto C^\ast(\Gamma)
\]

\[
(M,D)\mapsto \text{spectral triple}
\]

\[
A\mapsto K_\ast(A)
\]

\[
D\mapsto \operatorname{ind}(D).
\]



\section{H. Metric, Convex, and Optimization Networks}



\subsection{H1. Network metric--embedding--graphs--Banach--algorithms}

Nodes:

\[
\text{graphs}
\]

\[
\text{metric finite}
\]

\[
\text{Banach spaces}
\]

\[
\text{manifolds metric}
\]

\[
\text{expanders}
\]

\[
\text{optimal transport}
\]

\[
\text{metric embeddings}
\]

\[
\text{sketching}
\]

\[
\text{nearest neighbor search}
\]

\[
\text{communication complexity}.
\]

Edges:

\[
G\mapsto d_G
\]

\[
(X,d)\mapsto \ell_p^n
\]

\[
(X,d)\mapsto \operatorname{Lip}_1(X)
\]

\[
\mu,\nu\mapsto W_p(\mu,\nu)
\]

\[
\text{communication problem}\mapsto \text{matrix/metric}.
\]



\subsection{H2. Network convexity--SOS--moments--optimization--verification data}

Nodes:

\[
\text{semialgebraic sets}
\]

\[
\text{polynomial optimization}
\]

\[
\text{sum-of-squares}
\]

\[
\text{moment problems}
\]

\[
\text{convex bodies}
\]

\[
\text{linear programming}
\]

\[
\text{semidefinite programming}
\]

\[
\text{control theory}
\]

\[
\text{Lyapunov functions}
\]

\[
\text{combinatorial optimization}
\]

\[
\text{proof complexity}.
\]

Edges:

\[
S={f_i\ge 0}\mapsto \text{quadratic module}
\]

\[
\text{polynomial optimization}\mapsto \text{moment relaxation}
\]

\[
\text{polytope}\mapsto \text{slack matrix}
\]

\[
\text{dynamical system}\mapsto \text{Lyapunov verification data}
\]

\[
\text{combinatorial problem}\mapsto \text{LP/SDP hierarchy}.
\]



\subsection{H3. Network convex geometry--isoperimetry--concentration--probability}

Nodes:

\[
\text{convex bodies}
\]

\[
\text{Brunn--Minkowski theory}
\]

\[
\text{isoperimetric inequalities}
\]

\[
\text{log-concave measures}
\]

\[
\text{concentration of measure}
\]

\[
\text{transport inequalities}
\]

\[
\text{random polytopes}
\]

\[
\text{asymptotic geometric analysis}
\]

\[
\text{high-dimensional probability}.
\]

Edges:

\[
K\mapsto h_K
\]

\[
K\mapsto \operatorname{Vol}(K+tB)
\]

\[
\mu\mapsto \text{concentration profile}
\]

\[
K\mapsto \text{isotropic constant}
\]

\[
\mu,\nu\mapsto T_{\mu\to\nu}.
\]



\subsection{H4. Network tropical algebra--amoebas--asymptotic analysis}

Nodes:

\[
\text{complex varieties}
\]

\[
\text{amoebas}
\]

\[
\text{tropical varieties}
\]

\[
\text{Newton polytopes}
\]

\[
\text{large deviations}
\]

\[
\text{WKB asymptotics}
\]

\[
\text{idempotent analysis}
\]

\[
\text{max-plus algebra}
\]

\[
\text{optimization}
\]

\[
\text{control}.
\]

Edges:

\[
z\mapsto \log|z|
\]

\[
\varepsilon\log(e^{f/\varepsilon}+e^{g/\varepsilon})\mapsto \max(f,g)
\]

\[
\text{Hamilton--Jacobi}\mapsto \text{max-plus linearity}
\]

\[
\text{polynomial}\mapsto \text{Newton polytope}.
\]



\section{Composite Mega-networks}

Le networks precedenti diventano potentissime quando compongono.



\subsection{M1. Mega-network syntax--geometry--complexity}

\[
\begin{gathered}
\text{finite presentation} \to \text{category of models} \to \text{representation space} \to \text{invariant quotient} \\
\to \text{tropicalization} \to \text{matroid/fan} \to \text{CSP} \to \text{proof system}.
\end{gathered}
\]

Overall transfer:

\[
\begin{gathered}
\text{complexity of classificazione} \to \text{complexity of the invariants} \to \text{complexity tropical} \\
\to \text{complexity CSP} \to \text{complexity of prova}.
\end{gathered}
\]



\subsection{M2. Mega-network topology--quantum--operators--algorithms}

\[
\begin{gathered}
\text{varieties/knots} \to \text{fundamental groups} \to \text{representations} \to \text{operator algebras} \\
\to \text{TQFT} \to \text{quantum circuits} \to \text{tensor networks} \to \text{SDP relaxations}.
\end{gathered}
\]

Observables lungo the network:

\[
\begin{gathered}
\text{characters} \to \text{operator traces} \to \text{quantum amplitudes} \to \text{correlations} \\
\to \text{moments SDP}.
\end{gathered}
\]



\subsection{M3. Mega-network arithmetic--dynamics--languages--compression}

\[
\begin{gathered}
\text{number/arithmetic variety} \to \text{finite reductions} \to \text{sequences} \to \text{automata} \\
\to \text{languages} \to \text{compression} \to \text{complexity descrittiva}.
\end{gathered}
\]

Observables:

\[
\text{small primes}
\]

\[
\text{digital prefixes}
\]

\[
\text{factors of words}
\]

\[
\text{dimension of the automaton}
\]

\[
\text{parse size}
\]

\[
\text{entropy subword}.
\]

Question:

\[
\text{when does an arithmetic object produce linguistically simple observable profiles?}
\]



\subsection{M4. Mega-network universal local-to-global}

\[
\begin{gathered}
\text{object globale} \to \text{cover locale} \to \text{presheaf} \to \text{cohomology} \\
\to \text{obstruction class} \to \text{verification data} \to \text{algorithm/proof}.
\end{gathered}
\]

Principle:

\[
\boxed{
\text{every global failure can be presented as a gluing failure.}
}
\]



\subsection{M5. Mega-network deformation--verification datum--formalization}

\[
\begin{gathered}
\text{geometric problem} \to \text{dg-Lie controller} \to \text{obstruction tower} \to \text{finite verification data} \\
\to \text{formal proof} \to \text{machine-checkable object}.
\end{gathered}
\]

Observables:

\[
H^0,H^1,H^2,\ldots
\]

Verification data:

\[
\text{vanishing}
\]

\[
\text{obstruction nonzero}
\]

\[
\text{formal smoothness}
\]

\[
\text{rigidity}
\]

\[
\text{extension to order }n.
\]



\subsection{M6. Composite analytic mega-network}

Prima catena:

\[
\begin{gathered}
\text{functions/distributions} \to \text{Fourier/wavelet coefficients} \to \text{function spaces} \to \text{operators} \\
\to \text{PDE} \to \text{semigroups} \to \text{kernel} \to \text{spectra} \\
\to \text{geometry} \to \text{topology}.
\end{gathered}
\]

Seconda catena:

\[
\begin{gathered}
\text{PDE singolare} \to \text{microlocal data} \to \text{simplectic geometry} \to \text{sheaf category} \\
\to \text{Floer/Fukaya} \to \text{mirror algebra}.
\end{gathered}
\]

Terza catena:

\[
\begin{gathered}
\text{rumore} \to \text{random distributions} \to \text{SPDE} \to \text{regularity structures} \\
\to \text{Hopf algebras} \to \text{renormalization} \to \text{QFT}.
\end{gathered}
\]

Quarta catena:

\[
\begin{gathered}
\text{microstructure} \to \text{homogenized tensor} \to \text{effective PDE} \to \text{macroscopic observable} \\
\to \text{inverse design}.
\end{gathered}
\]

Quinta catena:

\[
\begin{gathered}
\text{dynamical system} \to \text{transfer operator} \to \text{spectrum} \to \text{zeta function} \\
\to \text{entropy/pressure} \to \text{large deviations}.
\end{gathered}
\]



\subsection{M7. Mega-network discrete--continuous--logical--numerical}

\[
\begin{gathered}
\text{finite graph} \to \text{metric measure approximation} \to \text{continuum space} \to \text{PDE operator} \\
\to \text{discretized operator} \to \text{numerical solution} \to \text{interval verification data} \to \text{formal proof}.
\end{gathered}
\]

Ciclo:

\[
\begin{gathered}
\text{discrete} \to \text{continuous} \to \text{discrete} \to \text{verification datum}.
\end{gathered}
\]



\subsection{M8. Mega-network arithmetic--analytic--spectral}

\[
\begin{gathered}
\text{arithmetic variety} \to \text{rappresentazione of Galois} \to L\text{-function} \to \text{automorphic representation} \\
\to \text{spectral theory} \to \text{trace formula} \to \text{geometric cycles}.
\end{gathered}
\]

Observables:

\[
\text{Frobenius traces}
\]

\[
\text{Euler factors}
\]

\[
\text{Fourier coefficients}
\]

\[
\text{eigenvalues}
\]

\[
\text{periods}
\]

\[
\text{intersection numbers}.
\]



\subsection{M9. Mega-network probability--PDE--optimization--learning}

\[
\begin{gathered}
\text{data distribution} \to \text{empirical measure} \to \text{gradient flow} \to \text{PDE mean-field} \\
\to \text{risk minimization} \to \text{generalization verification data} \to \text{algorithm}.
\end{gathered}
\]

Edges:

\[
\mu\to \operatorname{Ent}(\mu)
\]

\[
\mu_t\to \partial_t\mu_t=\nabla\cdot(\mu_t\nabla V)
\]

\[
\text{training}\to \text{mean-field PDE}
\]

\[
\text{PDE}\to \text{convergence rate}.
\]



\subsection{M10. Mega-network analytic local-to-global}

\[
\begin{gathered}
\text{local estimates} \to \text{patching} \to \text{global regularity} \to \text{compactness} \\
\to \text{existence} \to \text{verification data}.
\end{gathered}
\]

Nodes:

\[
\text{elliptic estimates}
\]

\[
\text{partitions of unity}
\]

\[
\text{Sobolev embeddings}
\]

\[
\text{compactness}
\]

\[
\text{weak convergence}
\]

\[
\text{defect measures}
\]

\[
\text{concentration compactness}.
\]



\subsection{M11. Mega-network of bounded observation}

Scheme universal:

\[
X\mapsto \mathcal O_{\le r}(X).
\]

Examples:

\[
u\mapsto \widehat u|_{|\xi|\le \Lambda}
\]

\[
M\mapsto {\lambda_1,\ldots,\lambda_N}
\]

\[
\mu\mapsto (m_\alpha)_{|\alpha|\le d}
\]

\[
T\mapsto \operatorname{Tr}(T^k)_{k\le N}
\]

\[
f\mapsto \text{samples }f(x_i)
\]

\[
P\mapsto \text{boundary response up to time }T
\]

\[
\text{graph}\mapsto \text{subgraph counts up to size }k
\]

\[
\text{structure}\mapsto \text{sentences of quantifier rank }\le q.
\]

Principle:

\[
\boxed{
\begin{gathered}
\text{every bounded observation induces a partition;}\\
\text{complexity is the geometry of those classes.}
\end{gathered}
}
\]



\section{Most Fertile Directions}

The strongest directions that emerge are the following.



\subsection{Universal sheafification of complexity problems}

Dato a problema globale (P), costruire:

\[
P\mapsto \mathcal F_P.
\]

Poi misurare:

\[
\text{minimo raggio locale or degree coomological che vede the fallimento.}
\]

Applicabile a:

\[
\text{CSP}
\]

\[
\text{database}
\]

\[
\text{distributed computing}
\]

\[
\text{proof complexity}
\]

\[
\text{matching}
\]

\[
\text{graph coloring}
\]

\[
\text{SAT}
\]

\[
\text{learning locale}
\]

\[
\text{bundle theory}.
\]



\subsection{Tensorization of non-tensorial structures}

Cercare the tensore nascosto:

\[
\text{groups}\to \text{tensore of commutatore}
\]

\[
\text{topological spaces}\to \text{cup-product tensor}
\]

\[
\text{PDE}\to \text{kernel tensor}
\]

\[
\text{database joins}\to \text{incidence tensor}
\]

\[
\text{proof systems}\to \text{coefficient tensor}
\]

\[
\text{quantum states}\to \text{entanglement tensor}
\]

\[
\text{statistical models}\to \text{probability tensor}.
\]

Poi usare:

\[
\text{rank}
\]

\[
\text{border rank}
\]

\[
\text{slice rank}
\]

\[
\text{flattening rank}
\]

per trasferire lower bounds.



\subsection{Fibres as primary objects}

Non studiare solo:

\[
I:X\to Y,
\]

ma:

\[
\operatorname{Fib}_I(y).
\]

Examples:

\[
\text{spaces with the same homology}
\]

\[
\text{graphs with same spectrum}
\]

\[
\text{modelli statistici with same moments}
\]

\[
\text{varieties with same tropicalization}
\]

\[
\text{groups with same finite quotients up to a }N
\]

\[
\text{functions with the same low coefficients}
\]

\[
\text{PDEs with the same boundary data}
\]

\[
\text{neural networks with the same function}.
\]



\subsection{Transfer of verification data}

Per every arco chiedere:

\[
\text{can target verification data be pulled back?}
\]

\[
\text{can source verification data be pushed forward?}
\]

verification data to study:

\[
\text{SOS}
\]

\[
\text{SDP}
\]

\[
\text{energy}
\]

\[
\text{index}
\]

\[
\text{cohomology}
\]

\[
\text{entropy}
\]

\[
\text{Lyapunov}
\]

\[
\text{barrier}
\]

\[
\text{interval enclosure}
\]

\[
\text{formal proof term}.
\]



\subsection{Budgeted Morita Theory}

Two theories are equivalent with budget:

\[
A\simeq_{\mathcal H}B.
\]

Possible classes of overhead:

\[
\text{linear}
\]

\[
\text{polynomial}
\]

\[
\text{quasi-polynomial}
\]

\[
\text{elementary}
\]

\[
\text{computable}
\]

\[
\text{non-computable}.
\]

This direction makes it possible to classify not only objects, but whole theories.



\subsection{Microlocal--sheaf--PDE--symplectic}

\[
\begin{gathered}
\text{PDE singularities} \to WF \to \text{Lagrangians} \to \text{microlocal sheaves} \\
\to \text{Fukaya categories}.
\end{gathered}
\]

Possible use:

\[
\text{analytic lower bounds}
\to
\text{categorical lower bounds}.
\]



\subsection{SPDE--regularity structures--Hopf algebras--QFT}

\[
\begin{gathered}
\text{noise} \to \text{singular PDE} \to \text{renormalization} \to \text{Hopf algebra} \\
\to \text{field theory}.
\end{gathered}
\]

Here the correct presentation is the renormalized object, not the raw datum.



\subsection{Homogenization--inverse design--optimization--learning}

\[
\begin{gathered}
\text{microstructure} \to \text{effective tensor} \to \text{target material} \to \text{inverse problem} \\
\to \text{generative design}.
\end{gathered}
\]

The fibre contains all equivalent microscopic materials.



\subsection{Computable analysis--proof mining--verified numerics}

\[
\begin{gathered}
\text{existence theorem} \to \text{effective modulus} \to \text{numerical scheme} \to \text{interval verification data} \\
\to \text{formal proof}.
\end{gathered}
\]

This is the network that turns qualitative analysis into verifiable mathematics.



\subsection{Dynamical systems--transfer operators--zeta--spectral geometry}

\[
\begin{gathered}
T \to \mathcal L_T \to \operatorname{Spec} \to \zeta_T \\
\to \text{periodic orbit data}.
\end{gathered}
\]

It is an ideal network for transporting dynamics into spectral analysis.



\subsection{Banach geometry--metric embeddings--algorithms}

\[
\begin{gathered}
\text{metric spaces} \to \text{Banach embeddings} \to \text{distortion} \to \text{sketching} \\
\to \text{communication complexity}.
\end{gathered}
\]

It is a natural network for algorithmic lower bounds.



\subsection{Riemann--Hilbert--D-modules--monodromy--asymptotics}

\[
\begin{gathered}
\text{PDE/ODE} \to D\text{-module} \to \text{perverse sheaf} \to \text{monodromy} \\
\to \text{Stokes data}.
\end{gathered}
\]

This turns differential analysis into constructible topology and algebra.



\section{Operational Sheet for Building a Transfer Package}

For each edge of the atlas one should fill in a sheet.



\subsection{Standard Sheet}

Dato:

\[
\mathfrak T:A\to B,
\]

specificare:

\subsection*{Source Objects}

\[
\mathcal X_A.
\]

\subsection*{Target Objects}

\[
\mathcal X_B.
\]

\subsection*{Source Descriptions}

\[
\mathcal D_A.
\]

\subsection*{Target Descriptions}

\[
\mathcal D_B.
\]

\subsection*{Compiler}

\[
\widehat F:\mathcal D_A\to\mathcal D_B.
\]

\subsection*{Presentational Overhead}

\[
\kappa_B(\widehat F(d))
\le
f(\kappa_A(d)).
\]

\subsection*{Target Observables}

\[
\mathfrak O_B.
\]

\subsection*{Observational Pullback}

\[
F^\ast:\mathfrak O_B\to\mathfrak O_A.
\]

\subsection*{Observational Overhead}

\[
\chi_A(F^\ast O)
\le
g(\chi_B(O)).
\]

\subsection*{Verification Data}

\[
V_A\leftrightarrow V_B.
\]

\subsection*{Fibres}

\[
F^{-1}(y).
\]

\subsection*{Moves in the Fibre}

\[
\text{generators, moves, deformations, homotopies, refinements}.
\]

\subsection*{Bridge profile}

\[
B_y(n)=\text{maximum/minimum cost for navigating in the fibre with budget }n.
\]

\subsection*{Measures}

\[
\mu_A\mapsto \mu_B.
\]

\subsection*{Limits}

\[
X_n\to X
\quad\mapsto\quad
F(X_n)\to F(X).
\]

\subsection*{Fertile Questions}

\[
\text{separazione}
\]

\[
\text{indistinguishability}
\]

\[
\text{normalizzazione}
\]

\[
\text{noncomputability}
\]

\[
\text{lower bounds}
\]

\[
\text{moduli}
\]

\[
\text{Verification Data}
\]

\[
\text{fibre}
\]

\[
\text{distortion}.
\]



\section{Final Global Map}

The previous linear picture is useful as a mnemonic, but it is too small to represent the atlas: it hides algebra, representation theory, graph and matroid methods, moduli, arithmetic geometry, and several discrete-continuous bridges. A better final map should be read as three complementary layers.

First, the large thematic regions are:

\[
\begin{gathered}
\text{logic, syntax, proof systems, formalization}\\
\to \text{algebra, representations, moduli, homological algebra}\\
\to \text{combinatorics, graphs, matroids, codes, discrete structures}\\
\to \text{geometry, topology, sheaves, persistence, categories}\\
\to \text{analysis, PDE, spectra, dynamics, probability}\\
\to \text{arithmetic, automata, languages, symbolic dynamics, computability}\\
\to \text{quantum, operator algebras, optimization, learning, numerics}.
\end{gathered}
\]

Second, the main access hubs are:

\[
\begin{array}{cccc}
\text{syntax} & \text{presentations} & \text{normal forms} & \text{compilers} \\
\text{observables} & \text{moments} & \text{rank profiles} & \text{spectra} \\
\text{rings/algebras} & \text{modules} & \text{representations} & \text{moduli} \\
\text{graphs} & \text{matroids} & \text{codes} & \text{finite structures} \\
\text{sheaves} & \text{cohomology} & \text{obstructions} & \text{local-to-global data} \\
\text{operators} & \text{kernels} & \text{semigroups} & \text{resolvents} \\
\text{probability} & \text{entropy} & \text{large deviations} & \text{optimal transport} \\
\text{numerics} & \text{approximations} & \text{verification data} & \text{formal proofs}.
\end{array}
\]

Third, the atlas is organized by bridges rather than by isolated subjects:

\[
\begin{gathered}
\text{syntax} \leftrightarrow \text{semantics} \leftrightarrow \text{CSP/database/proof complexity},\\
\text{algebra} \leftrightarrow \text{representations} \leftrightarrow \text{moduli} \leftrightarrow \text{invariants},\\
\text{combinatorics} \leftrightarrow \text{graphs} \leftrightarrow \text{matroids} \leftrightarrow \text{extremal/counting methods},\\
\text{linear algebra} \leftrightarrow \text{matroids} \leftrightarrow \text{codes} \leftrightarrow \text{rank methods},\\
\text{graphs} \leftrightarrow \text{limits} \leftrightarrow \text{operators} \leftrightarrow \text{spectra},\\
\text{geometry} \leftrightarrow \text{sheaves} \leftrightarrow \text{cohomology} \leftrightarrow \text{obstructions},\\
\text{PDE} \leftrightarrow \text{semigroups} \leftrightarrow \text{kernels} \leftrightarrow \text{probability},\\
\text{arithmetic} \leftrightarrow \text{Galois/automata/languages} \leftrightarrow \text{compression},\\
\text{optimization} \leftrightarrow \text{convexity/SOS/moments} \leftrightarrow \text{verification data}.
\end{gathered}
\]

This map is still not a theorem. It is a navigation chart: every arrow represents a family of possible transfer packages, and every useful package must specify the compiler, observable pullback, verification data, fibre geometry, and overhead functions involved.

\section{Final Conceptual Summary}

The atlas can be summarized in seven layers.

\subsection{Layer 1: objects}

\[
\begin{gathered}
\text{functions, spaces, groups, varieties, operators, measures, categories},\\
\text{graphs, models, processes, dynamical systems, PDEs},\\
\text{logical theories, neural networks, quantum objects}.
\end{gathered}
\]

\subsection{Layer 2: presentations}

\[
\begin{gathered}
\text{equations, generators and relations, coefficients, kernels},\\
\text{triangulations, local data, graphs, matrices, tensors},\\
\text{programs, networks, schemes, samples}.
\end{gathered}
\]

\subsection{Layer 3: observables}

\[
\begin{gathered}
\text{spectra, moments, entropies, cohomology, homology},\\
\text{samples, queries, correlations, invariants, heat kernels},\\
\text{traces, barcodes, ranks, dimensions, growth profiles}.
\end{gathered}
\]

\subsection{Layer 4: verification data}

\[
\begin{gathered}
\text{SOS and SDP data, energy bounds, index data, convex dual data},\\
\text{cohomological obstructions, a priori estimates, stability data},\\
\text{entropy data, Lyapunov data, barriers, interval enclosures},\\
\text{proof terms and formal proof data}.
\end{gathered}
\]

\subsection{Layer 5: fibres}

\[
\begin{gathered}
\text{objects indistinguishable under bounded observables},\\
\text{equivalent descriptions and non-identifiable parameters},\\
\text{microstructures with the same limit},\\
\text{functions with the same samples},\\
\text{operators with the same truncated spectrum},\\
\text{graphs with the same small subgraphs},\\
\text{theories with the same finite models},\\
\text{neural networks with the same function},\\
\text{PDEs with the same boundary data}.
\end{gathered}
\]

\subsection{Layer 6: scales}

\[
\begin{gathered}
\text{cutoff, precision, time, frequency, mesh, quantifier rank},\\
\text{dimension, degree, number of samples, number of states},\\
\text{SDP level, perturbative order, semiclassical scale},\\
\text{microscopic and macroscopic scales}.
\end{gathered}
\]

\subsection{Layer 7: transfer}

\[
\begin{gathered}
\text{compilers between descriptions, pullbacks of observables},\\
\text{pushforwards and pullbacks of verification data},\\
\text{transfer of fibres, limits, measures, regularity, scales},\\
\text{transfer of lower bounds}.
\end{gathered}
\]

\section{Final Formula of the Atlas}

The whole project can be condensed as follows:

\[
\boxed{
\text{every mathematical theory has many observable profiles;}
}
\]

\[
\boxed{
\text{every observable profile is produced by a system of observables;}
}
\]

\[
\boxed{
\text{every system of observables induces fibres;}
}
\]

\[
\boxed{
\text{every fibre contains hidden information;}
}
\]

\[
\boxed{
\begin{gathered}
\text{every transfer package measures the cost of passing between}\\
\text{observable profiles, descriptions, verification data, and fibres;}
\end{gathered}
}
\]

\[
\boxed{
\text{complexity is the quantitative geometry of access to objects.}
}
\]

In an even shorter form:

\[
\boxed{
\begin{gathered}
\text{Presentation Theory + Transfer Packages}\\
\Longleftrightarrow\\
\text{global geometry of mathematical observable profiles}\\
\text{and of their translation costs.}
\end{gathered}
}
\]

This atlas collects a network of networks connecting logic, algebra, combinatorics, geometry, analysis, probability, topology, arithmetic, algorithms, mathematical physics, optimization, learning, numerics, and formalization through a single language: \textbf{descriptions, observables, verification data, fibres, scales, and controlled transfers}.



\begin{thebibliography}{99}
\bibitem{presentation-theory}
L. Blanchi, \emph{Presentation Theory I: Description Systems, Costs, and Observable Fibres}, AI-assisted math note, 2026.

\bibitem{presentation-theory-ii}
L. Blanchi, \emph{Presentation Theory II: Controlled Transfer, Observable Budgets, and the Geometry of Fibres}, AI-assisted math note, 2026.

\bibitem{gromov}
M. Gromov, \emph{Metric Structures for Riemannian and Non-Riemannian Spaces}, Birkhauser, 1999.

\bibitem{lawvere}
F. W. Lawvere and S. H. Schanuel, \emph{Conceptual Mathematics: A First Introduction to Categories}, Cambridge University Press, 2009.
\end{thebibliography}

\end{document}
